\documentclass[11pt]{article}

\usepackage[T1]{fontenc}
\usepackage{lmodern}
\usepackage{amsmath, amssymb, amsthm}
\usepackage{boldtensors}
\usepackage{booktabs}
\usepackage{microtype}
\usepackage[margin=1.15in]{geometry}
\usepackage[colorlinks, linkcolor=blue, citecolor=blue]{hyperref}
\usepackage[numbers, sort&compress]{natbib}
\setlength{\emergencystretch}{2.5em}

\newcommand{\F}{~F}
\newcommand{\C}{~C}
\newcommand{\Id}{~I}
\renewcommand{\P}{~P}
\newcommand{\eps}{~\epsilon}
\newcommand{\edev}{~e}
\newcommand{\sdev}{~s}
\newcommand{\sdevx}{~s_{\!x}}
\newcommand{\btau}{~\tau}
\newcommand{\disp}{~\varphi}
\newcommand{\testd}{~\psi}
\newcommand{\bx}{~x}
\newcommand{\bX}{~X}
\newcommand{\bT}{~T}
\newcommand{\bB}{~B}
\newcommand{\bN}{~N}
\newcommand{\Syms}{\mathbb{S}}
\newcommand{\Reals}{\mathbb{R}}
\newcommand{\Grad}{\operatorname{Grad}}
\newcommand{\Div}{\operatorname{Div}}
\newcommand{\dev}{\operatorname{dev}}
\newcommand{\trace}{\operatorname{tr}}
\newcommand{\sym}{\operatorname{sym}}
\newcommand{\transpose}[1]{#1^{\mathsf{T}}}
\newcommand{\norm}[1]{\left\lVert #1 \right\rVert}
\newcommand{\diff}{\,\mathrm{d}}
\newcommand{\supvol}{^{\mathrm{vol}}}
\newcommand{\supdev}{^{\mathrm{dev}}}
\newcommand{\supiso}{^{\mathrm{iso}}}
\newcommand{\code}[1]{\texttt{#1}}

\theoremstyle{definition}
\newtheorem{remark}{Remark}
\theoremstyle{plain}
\newtheorem{proposition}{Proposition}

\title{\textbf{HW3L}\\[0.3em]
  \large A Three-Field Hu--Washizu Tetrahedron in Logarithmic Strain\\[0.2em]
  \normalsize Research note}
\author{Alejandro Mota}
\date{\today}

\begin{document}
\maketitle

\begin{abstract}
\noindent
HW3L (Hu--Washizu, three fields, logarithmic strain) is a tetrahedral element
for large-deformation deviatoric (trace-free, shape-changing) plasticity, aimed at removing the
stabilization that two families of methods carry: the composite tetrahedron,
a tetrahedron subdivided into four linear tetrahedra with one element-constant
pressure, which relieves volumetric locking and introduces soft modes that a
penalty energy removes; and variational-multiscale methods, which add a
subgrid pressure proportional to the time step that reaches the plastic return
map, the algorithm that projects a trial stress onto the yield surface. Three fields --- the motion, a volumetric strain and a pressure --- are
coupled through a Hu--Washizu functional, in which the two auxiliary fields are
independent of the motion and tied to it by a constraint term. The deviatoric
response is evaluated pointwise at quadrature, so the constitutive law is an
unmodified small-strain algorithm, and both auxiliary fields are discontinuous
across elements, so their elimination is element-local. The stored energy must
split exactly into volumetric and isochoric (volume-preserving) parts, which
five of the seven materials implemented in the finite-element code Norma do.
With both auxiliary fields in one discrete space the functional reduces
exactly to mean dilatation, the replacement of the volumetric strain by its
element-wise projection, and what is open is the choice of pair: the discrete
space for the displacement together with the discrete space for the pressure.
Volumetric locking and spurious soft modes are failures of the two hypotheses
of Brezzi's theorem for that pair: the inf-sup condition, that every discrete
pressure is resisted by some discrete displacement with a constant independent
of the mesh, and coercivity of the deviatoric form on the discretely isochoric
subspace, that no discrete displacement the pressure does not resist has
vanishing deviatoric energy. The first is settled by measurement over a mesh
sequence against a control pair of established stability: of seven
element-local pairs two are stable, the three-dimensional Crouzeix--Raviart
pair --- quadratic displacement enriched with face and interior bubbles,
polynomials on one element that vanish on a prescribed part of its boundary,
over a linear discontinuous pressure --- at a constant of $0.2968$, and the
constant-pressure pair with face bubbles, whose constant decreases toward a
limit near $0.35$; the five others decay. The face bubbles are global, and the
Crouzeix--Raviart pair has $2.6$ times the degrees of freedom of the composite
tetrahedron. The second holds in elasticity for every conforming pair, a pair
whose displacement space is continuous across elements, and under a plastic
tangent that removes the fraction $\eta$ of the
shear stiffness along the flow direction ($\eta = 1$ for perfect plasticity)
an identity bounds its constant below by $(1-\eta)/2$ for every pair, so it
cannot separate them. The claim of this note is that the pairs differ in
whether their soft modes are isochoric, and it is supported on the linearized
operator: under a plastic tangent the softest modes of a constant-pressure
element carry a dilatation of $0.20$ to $0.27$ that does not decay under
refinement, those of the Crouzeix--Raviart pair one that decays as the square
of the mesh size $h$, and those of the Taylor--Hood pair (continuous quadratic
displacement over continuous linear pressure) are dilatation almost entirely;
in a confined plastic zone with a varying flow direction the Crouzeix--Raviart
pair has no spurious mode among its twenty softest at the finest mesh, and,
with the whole boundary fixed, the constant-pressure pairs retain between five
and fifteen. On deformed
configurations the inf-sup constant decreases with the shear of the
deformation and by $7\%$ under a volume change of $2.3$ that shears little.
A fourteen-point
degree-five quadrature rule and a nodal fifteen-node basis are sufficient, and
the explicit stable step is $3.3$ times smaller than a linear tetrahedron's.
The smoothness of the equivalent plastic strain, the accumulated norm of the
plastic strain increments, in a nonlinear solve is not measured here.
\end{abstract}

\section{Problem statement}
\label{sec:problem}

Tetrahedra are the only practical choice for meshing complex geometry
automatically. Their accuracy is insufficient in four regimes: near-incompressible
response at bulk-to-shear modulus ratios $\kappa/\mu \lesssim 10^4$, with
$\kappa$ the bulk modulus and $\mu$ the shear modulus; isochoric
(volume-preserving) plasticity; large rotation; and weak shocks, pressure
waves of small amplitude, in materials of high $\kappa/\mu$. The target is
one spatial formulation for all four, with the time integrator the only
difference between the quasi-static and the dynamic regimes.

Two established remedies fail, and each failure is the failure of one
hypothesis of the same theorem; this section identifies the hypothesis in
each case.

\paragraph{Constant pressure relieves locking and introduces soft modes.}
The composite tetrahedron, a tetrahedron subdivided into four linear
tetrahedra with one element-constant pressure
\citep{Thoutireddy.etal:2002, Ostien.etal:2016}, and its five-field extension
\citep{Foulk.etal:2021} relieve volumetric locking --- the stiffening of a
discretization whose displacement space contains too few discrete isochoric
fields --- with that pressure. \citet{Foulk.etal:2021} report that this
choice --- and it is the choice, not the element --- produces a deleterious
soft mode ``common to both the quadratic and composite tetrahedral element with
constant pressure formulations,'' remedied there by a convex energy penalty.
The mechanism is a coercivity failure on the discretely isochoric subspace: an
element-constant multiplier does not resist displacement patterns whose volume
change varies within the element and averages to zero, so those patterns carry
only shear energy. In elasticity every projected-pressure element has such
patterns, and a stable pair keeps their energy bounded away from zero; in a
confined plastic zone the shear response softens along the flow direction and
the pattern can approach zero energy. \S\ref{sec:assembled} shows that on the
quadratic tetrahedron these patterns are not exactly of zero energy, and
\S\ref{sec:prediction} states the question that count does not settle.

\paragraph{Residual-based stabilization contaminates the deviator.}
Variational-multiscale stabilization of linear tetrahedra
\citep{Scovazzi:2015} introduces a subgrid pressure of the form
\begin{equation}
  \label{eq:vms}
  p' = c_1 \, \tau_{\mathrm s} \left( \kappa \Div \dot{\disp} - \dot{p} \right),
  \qquad
  \tau_{\mathrm s} = c_2 \, \Delta t,
\end{equation}
with $c_1$ and $c_2$ tuned coefficients, which is added to the stress before
the constitutive update is evaluated. Two consequences follow. The spatial
solution depends on the time step $\Delta t$ through $\tau_{\mathrm s}$, so
there is no quasi-static limit independent of it. The coefficient $c_2$ scales
a term that reaches the deviator, and therefore the plastic return map --- the
algorithm that projects a trial stress onto the yield surface --- which is the
origin of the oscillations in the equivalent plastic strain, the accumulated
norm of the plastic strain increments. No choice of $c_2$ removes this,
because the term enters the deviator by construction.

\paragraph{The common cause.} Brezzi's theorem for the mixed problem asks two things of
the pair --- the discrete space $U_h$ for the displacement together with the
discrete space $Q_h$ for the pressure
\citep[\S3.4, \S5.2]{BoffiBrezziFortin:2013}: the inf-sup condition,
that every pressure mode be resisted by some displacement, whose failure is
locking on one side and spurious pressure on the other; and coercivity of the
deviatoric form on the discretely isochoric subspace $\ker ~G_h$, that no
displacement the pressure does not resist have zero deviatoric energy, whose
failure is a soft mode.
The first is a property of the pair alone and is settled in \S\ref{sec:beta}.
The second depends on the material tangent as well as on the pair, and
\S\ref{sec:prediction} states what this note claims about it. Neither
difficulty is a stabilization problem: the remedy for the first is a pair
satisfying the inf-sup condition uniformly in the mesh size, and the second
must then be checked where the material can violate it.

\section{Continuum formulation}

This section states the kinematics, the stored energy, the three-field
functional, its reduction to a displacement functional, the requirements it
places on the material, and the stress it produces.

\subsection{Kinematics and the volumetric--isochoric split}
\label{sec:kinematics}

A body $\Omega \subset \Reals^3$ undergoes a motion $\bx = \disp(\bX, t)$ with
deformation gradient $\F := \Grad\disp$ and Jacobian $J := \det\F > 0$. The
multiplicative factorization
\begin{equation}
  \label{eq:isosplit}
  \F = J^{1/3}\bar{\F},
  \qquad
  \bar{\F} := J^{-1/3}\F,
  \qquad
  \det\bar{\F} = 1,
\end{equation}
separates the two kinematic contributions this formulation treats differently:
the mixed treatment acts on $J$ alone, and the plasticity on $\bar{\F}$ alone.

The mixed field is a scalar function of $J$ alone. Let
\begin{equation}
  \label{eq:thetadef}
  \theta = \theta(J),
  \qquad \theta'(J) > 0,
  \qquad \theta(1) = 0,
  \qquad \theta'(1) = 1,
\end{equation}
be a strictly increasing volumetric strain measure, normalized to vanish with
unit slope in the undeformed state. The normalization has two consequences. It
makes $\kappa$ in \eqref{eq:energy} the physical bulk modulus for every
admissible $\theta$, so that two instances of the formulation are calibrated to
the same material rather than to two different ones. And since
$\partial J/\partial\F = J\F^{-\mathsf T}$,
\begin{equation}
  \label{eq:dtheta}
  \frac{\partial\theta}{\partial\F} = \theta'(J)\, J\, \F^{-\mathsf T},
\end{equation}
so \eqref{eq:thetadef} also forces $\partial\theta/\partial\F = \Id$ at
$\F = \Id$ for every admissible $\theta$ (\S\ref{sec:fe-spaces}).

Two choices are used below and are developed in \S\ref{sec:instances}: the
logarithmic measure $\theta = \log J$, and $\theta = J - 1$. Both satisfy
\eqref{eq:thetadef}, and by Remark~\ref{rem:thetachoice} they nevertheless give
different elements.

\subsection{Stored energy}
\label{sec:energy}

The formulation is developed for a stored energy that splits exactly along
\eqref{eq:isosplit}, with a volumetric part quadratic in $\theta$:
\begin{equation}
  \label{eq:energy}
  W(\F) = W\supvol(\theta) + W\supiso(\bar{\F}),
  \qquad
  W\supvol(\theta) = \tfrac{\kappa}{2}\theta^2 .
\end{equation}
\S\ref{sec:matreq} states what each half of \eqref{eq:energy} requires of the
material.

Given \eqref{eq:energy}, the volumetric constitutive relation is linear,
\begin{equation}
  \label{eq:constitutive}
  p = \partial W\supvol/\partial\theta = \kappa\theta ,
\end{equation}
and it is this linearity --- not the choice of $\theta$, and not the
logarithmic measure specifically --- that makes the discrete pressure equal
$\kappa\bar\theta$ pointwise in \S\ref{sec:local}. The elimination itself rests
on less, as \S\ref{sec:matreq} shows.

The isochoric response is the $W\supiso$ the material supplies. It is
evaluated at the pointwise $\bar{\F}$ and is never projected, which is the
property \S\ref{sec:plasticity} requires. Plasticity acts entirely inside
$W\supiso$, through an isochoric internal variable evolved by a return map;
neither \eqref{eq:energy} nor \eqref{eq:constitutive} changes when it does.

\subsection{Three-field functional}
\label{sec:hw}

Only the volumetric response is treated mixedly. Introduce an independent
volumetric strain $\bar\theta$ and constrain it to $\theta(\disp)$ through a
multiplier $\bar p$, which at stationarity is the pressure:
\begin{equation}
  \label{eq:hw}
  \begin{split}
    \Pi[\disp, \bar\theta, \bar p] & :=
    \int_\Omega W\supvol(\bar\theta) \diff V
    + \int_\Omega W\supiso\!\left(\bar{\F}(\disp)\right) \diff V
    + \int_\Omega \bar p \left(\theta(\disp) - \bar\theta\right) \diff V
    \\
    & \quad
    - \int_\Omega \rho_0\, ~b \cdot \disp \diff V
    - \int_{\partial_T\Omega} \bT \cdot \disp \diff S ,
  \end{split}
\end{equation}
with $\rho_0$ the reference density, $~b$ the body force per unit mass, and
$\bT$ the traction prescribed on the part $\partial_T\Omega$ of the boundary.
The isochoric energy is evaluated at the pointwise $\bar{\F}(\disp)$
--- for the Hencky instance of \S\ref{sec:instances}, at $\edev(\disp)$ ---
and carries no auxiliary field; Remark~\ref{rem:whynodev} gives the reason.

\subsection{Reduction to the mean-dilatation formulation}
\label{sec:meandilatation}

With $\bar\theta$ and $\bar p$ in one discrete space, \eqref{eq:hw} reduces to
the mean-dilatation formulation of \citet{SimoTaylorPister:1985}; this section
gives the reduction.

Put $\bar\theta$ and $\bar p$ in the same discrete space $Q_h$, as
\S\ref{sec:fe-spaces} requires. Stationarity in $\bar\theta$
then gives $\bar p = \kappa\bar\theta$ exactly rather than after a
projection, and stationarity in $\bar p$ gives $\bar\theta = \mathcal{P}_h\,\theta$,
the $L^2$ projection of $\theta$ onto $Q_h$. The constraint term of
\eqref{eq:hw} then vanishes identically:
\begin{equation}
  \label{eq:constraint-vanishes}
  \int_\Omega \bar p \left(\theta - \bar\theta\right) \diff V
  = \kappa \int_\Omega \bar\theta \left(\theta - \mathcal{P}_h \theta\right) \diff V = 0,
\end{equation}
because $\kappa\bar\theta \in Q_h$ while $(\theta - \mathcal{P}_h\theta) \perp Q_h$ by
the definition of the projection. What is left is a functional in the motion
alone,
\begin{equation}
  \label{eq:collapsed}
  \Pi[\disp] = \int_\Omega W\supiso\!\left(\bar{\F}(\disp)\right) \diff V
             + \int_\Omega \tfrac{\kappa}{2}\left(\mathcal{P}_h \theta\right)^2 \diff V,
\end{equation}
whose stationarity conditions are exactly the three-field equations of
\S\ref{sec:variations}.

Equation~\eqref{eq:collapsed} is the mean-dilatation formulation of
\citet{SimoTaylorPister:1985}, in which the volumetric strain is replaced by
its projection onto an element-wise space, and the mechanism of the $\bar B$
family (projected strain--displacement operator) and the $\bar F$ family
(projected deformation gradient) of elements; \citet[\S8.12]{BoffiBrezziFortin:2013}
set out the relation between reduced integration, projected pressure, and the
mixed method it is equivalent to. The reduction is exact by
\eqref{eq:constraint-vanishes}, and no property of $\theta$ or of
$W\supvol$ enters it: the exactness is a property of the discretization, that
the two auxiliary fields share one space. What the material must supply is
stated in \S\ref{sec:matreq}, and the quantity the formulation leaves open is
$Q_h$, which \S\ref{sec:fe-spaces} selects and \S\ref{sec:results} measures.

\subsection{Requirements on the material}
\label{sec:matreq}

This section states what each half of \eqref{eq:energy} requires of the
material, together with a third condition that is not a requirement of this
section, and reports which materials meet them.

\paragraph{An exact volumetric--isochoric split.} This condition is
structural. The
element evaluates $W$ at the pointwise deformation, subtracts the pointwise
volumetric energy and adds it back at the projected strain, so the subtraction
must leave no volumetric remainder:
\begin{equation}
  \label{eq:matsplit}
  W(\F) = W\supvol(J) + W\supiso(\bar{\F}),
  \qquad
  \bar{\F} := J^{-1/3}\F,
  \qquad
  \det\bar{\F} = 1 .
\end{equation}
Whether a material satisfies \eqref{eq:matsplit} is decidable without knowing
$W\supvol$: if the split exists then $W(\F) - W(\bar{\F}) = W\supvol(J) -
W\supvol(1)$ depends on $J$ alone, so holding $J$ fixed while varying the
isochoric part must leave the difference unmoved. Column~2 of
Table~\ref{tab:split} is that test, run over the seven materials implemented
in Norma, the finite-element code in which the element-level prototype of the
formulation is written; Carina, the code in which the formulation is to be
implemented, implements two of them.
Five of the seven satisfy it; the element-level implementation in Norma
admits only Hencky.

\paragraph{A quadratic $W\supvol$.} This condition is a convenience and not a
necessity. Stationarity in $\bar p$ gives $\bar\theta = \mathcal{P}_h\theta$ regardless of
the material, and stationarity in $\bar\theta$ then gives
\begin{equation}
  \label{eq:pbar-general}
  \bar p = \mathcal{P}_h\!\left(\partial W\supvol/\partial\bar\theta\right)
\end{equation}
in general --- the projection of a known function of the already
computed $\bar\theta$, one more mass-matrix solve with a nonlinear integrand
and nothing implicit. When $W\supvol$ is quadratic its derivative is linear,
$\kappa\bar\theta$ already lies in $Q_h$, the projection in
\eqref{eq:pbar-general} is the identity, and $\bar p = \kappa\bar\theta$ holds
pointwise; otherwise $\bar p = \mathcal{P}_h W\supvol{}'(\mathcal{P}_h\theta)$ is a
doubly projected pressure. The reduction \eqref{eq:collapsed} is exact either
way, with $W\supvol(\mathcal{P}_h\theta)$ in place of the quadratic, and the
elimination of \S\ref{sec:local} is linear either way; what the quadratic case
gives is the pointwise relation and one projection fewer.

\paragraph{The logarithmic measure itself.} This condition is required by
nothing in the present section. It is required by \S\ref{sec:plasticity}.

\begin{table}[tbp]
  \centering
  \caption{What the materials implemented in Norma satisfy. Column~2 is the
  spread of $W(\F) - W(\bar{\F})$ across twenty random isochoric parts at each
  of four fixed $J \in [0.85, 1.2]$, relative to its own magnitude; a spread
  below $10^{-12}$ means the split \eqref{eq:matsplit} holds to rounding
  error. Columns~3 and~4 are the
  relative residual of a least-squares fit of the extracted $W\supvol(J) =
  W(J^{1/3}\Id) - W(\Id)$ over $J \in [0.8, 1.25]$ to a quadratic in each
  candidate volumetric variable; in both exact cases the fitted coefficient is
  $\kappa/2$ to five digits. Two materials admit a quadratic $W\supvol$, and
  they disagree on the variable. Script: \code{prototype/materials.jl}, which
  needs Norma.}
  \label{tab:split}
  \begin{tabular}{lccc}
    \toprule
    & split \eqref{eq:matsplit} & \multicolumn{2}{c}{$W\supvol$ quadratic in} \\
    \cmidrule(lr){3-4}
    material & spread at fixed $J$ & $J - 1$ & $\log J$ \\
    \midrule
    Linear elastic              & $2.0\times10^{0}$   & ---                 & ---                 \\
    Saint-Venant Kirchhoff      & $3.0\times10^{1}$   & ---                 & ---                 \\
    Neo-Hookean                 & $4.2\times10^{-13}$ & $5.9\times10^{-2}$  & $1.3\times10^{-1}$  \\
    Reciprocal neo-Hookean      & $4.8\times10^{-13}$ & $1.8\times10^{-1}$  & $6.5\times10^{-3}$  \\
    Seth--Hill ($m = n = 1$)    & $4.3\times10^{-13}$ & $1.8\times10^{-1}$  & $6.5\times10^{-3}$  \\
    Hencky                      & $1.4\times10^{-13}$ & $1.8\times10^{-1}$  & $\mathbf{1.7\times10^{-15}}$ \\
    $J_2$, Simo--Hughes         & $5.9\times10^{-13}$ & $\mathbf{5.3\times10^{-15}}$ & $1.8\times10^{-1}$  \\
    \bottomrule
  \end{tabular}
\end{table}

\paragraph{The case developed here.} The formulation is developed in detail
for materials satisfying \eqref{eq:matsplit} with a quadratic
$W\supvol$, because that is where the pointwise relation holds and what both
implemented materials satisfy. Table~\ref{tab:split} admits exactly two, and they do
not agree on the variable: Hencky is quadratic in $\log J$ to
$1.7\times10^{-15}$, and the Simo--Hughes $J_2$ model that Norma and Carina
both implement is quadratic in $J-1$ to $5.3\times10^{-15}$. Each is exact in its
own variable and wrong by $18\%$ in the other, so material and volumetric
variable are not independently selectable, and neither choice generalizes to
the other material. \S\ref{sec:instances} develops both.

The three remaining split materials --- neo-Hookean, its reciprocal, and
Seth--Hill --- are admissible with one projection more
(Remark~\ref{rem:nonquadratic}) and are not developed here. What the
formulation requires of a material is the split, and five of the seven have
it.

\begin{remark}[Extending to a non-quadratic $W\supvol$]
\label{rem:nonquadratic}
Neither \S\ref{sec:hw} nor \S\ref{sec:meandilatation} is affected, and no
local Newton iteration is introduced. Two things change.

First, \eqref{eq:projection} acquires a second projection: $\bar\theta =
\mathcal{P}_h\theta$ from one mass-matrix solve as before, then $\bar p = \mathcal{P}_h
W\supvol{}'(\bar\theta)$ from a second solve with the same factorization and a
nonlinear integrand. The tangent contribution is
$\int W\supvol{}''(\bar\theta)\,\mathcal{P}_h(D\theta[\delta\disp])\,
\mathcal{P}_h(D\theta[\testd])\diff V$, a weighted integral against the same constant
mass matrix, and the factorization is reused across Newton iterations exactly
as in the quadratic case.

Second, the choice of volumetric variable is not determined by the material.
For a quadratic $W\supvol$ the variable is whichever one makes it quadratic,
and Table~\ref{tab:split} settles it. For a non-quadratic $W\supvol$ any
monotone $\theta(J)$ satisfying \eqref{eq:thetadef} is admissible, and by
Remark~\ref{rem:thetachoice} the choice of $\theta$ changes the discrete
element.
\end{remark}

\begin{remark}[The volumetric variable is a modeling choice]
\label{rem:thetachoice}
Projecting $\log J$ and projecting $J - 1$ do not give the same element.
Measured on one tetrahedron under a displacement quadratic in $\bX$, with the
element mean as the projection and a Hencky material so that both variants use
the same exact $W\supvol$, the element energies differ by
$9.4\times10^{-4}$ at $0.5\%$ nominal strain,
$1.0\times10^{-2}$ at $5\%$,
$5.6\times10^{-2}$ at $20\%$ and
$2.7\times10^{-1}$ at $50\%$
(\code{prototype/materials.jl}). The difference is $O(\text{strain})$ and
vanishes with it, so both are consistent discretizations of the same continuum
problem; they are different elements, and they differ by $5.6\%$ to $27\%$ at
$20\%$ to $50\%$ strain, the finite-deformation regime of
\S\ref{sec:problem}. Which variable is projected is therefore a modeling
decision.
\end{remark}

\subsection{The two admissible instances}
\label{sec:instances}

Table~\ref{tab:split} leaves two materials. They agree on everything
\S\ref{sec:kinematics}--\S\ref{sec:meandilatation} requires;
Table~\ref{tab:instances} lists where they differ.

\begin{table}[tbp]
  \centering
  \caption{The two instances of \eqref{eq:energy}. Everything above the rule is
  fixed by the generic development; everything below it is where the instances
  differ. $~b^{\mathrm e} := \F^{\mathrm e}\transpose{(\F^{\mathrm e})}$
  is the elastic left Cauchy--Green tensor, with $\F^{\mathrm e}$ the elastic
  part of the multiplicative decomposition $\F = \F^{\mathrm e}\F^{\mathrm p}$, and
  $\bar{~b}^{\mathrm e} := J^{-2/3}~b^{\mathrm e}$ its isochoric part.}
  \label{tab:instances}
  \begin{tabular}{lll}
    \toprule
                                & Hencky                       & Simo--Hughes \\
    \midrule
    volumetric variable $\theta$          & $\log J$           & $J - 1$ \\
    $\theta'(J)$                          & $1/J$              & $1$ \\
    $\partial\theta/\partial\F$           & $\F^{-\mathsf T}$  & $J\,\F^{-\mathsf T}$ \\
    $W\supvol(\theta)$                    & $\tfrac{\kappa}{2}(\log J)^2$ & $\tfrac{\kappa}{2}(J-1)^2$ \\
    \midrule
    isochoric argument                    & $\edev = \dev\eps$, $\eps = \tfrac12\log\C$
                                                               & $\bar{~b}^{\mathrm e}$ \\
    $W\supiso$                            & $\mu\norm{\edev}^2$
                                                               & $\tfrac{\mu}{2}(\trace\bar{~b}^{\mathrm e} - 3)$ \\
    return map                            & radial return in $\eps$ & radial return on $\bar{~b}^{\mathrm e}$ \\
    small-strain models usable unmodified & yes                & no \\
    implemented in Norma and Carina       & elasticity only    & yes, with consistent tangent \\
    \bottomrule
  \end{tabular}
\end{table}

The Hencky instance is the one \eqref{eq:hencky} defines: with
\begin{equation}
  \label{eq:hencky}
  \eps := \tfrac{1}{2}\log\C \in \Syms,
  \qquad \C := \transpose{\F}\F,
\end{equation}
evaluated by inverse scaling and squaring with Pad\'e approximants --- repeated
square roots of $\C$ followed by a rational approximant of the logarithm
\citep{AlMohy.Higham.Relton:2013} --- the
volumetric--deviatoric decomposition
\begin{equation}
  \label{eq:split}
  \eps = \tfrac{1}{3}\theta\Id + \edev,
  \qquad \theta = \trace\eps = \log J,
  \qquad \edev = \dev\eps,
\end{equation}
is orthogonal in the Frobenius inner product, and the split of
\eqref{eq:energy} is inherited from the strain measure rather than imposed on
it. The Simo--Hughes instance imposes it instead, through the $J^{-2/3}$ factor
in $\bar{~b}^{\mathrm e}$; $\det\bar{~b}^{\mathrm e} = 1$ then makes
$W\supiso$ independent of $J$ by construction. Both are exact, and
Table~\ref{tab:split} confirms both to machine precision. Three consequences
follow.

\paragraph{The generic development covers both.} Nothing in
\S\ref{sec:kinematics}--\S\ref{sec:meandilatation}, and nothing in the
discretization of \S\ref{sec:fe-spaces}--\S\ref{sec:local}, refers to either
column of Table~\ref{tab:instances}. The reduction \eqref{eq:collapsed}, the
pointwise relation $\bar p = \kappa\bar\theta$, and the element-local
elimination hold for any $\theta$ satisfying \eqref{eq:thetadef} and any
$W\supiso$ satisfying \eqref{eq:energy}.

\paragraph{The pairing at $\F = \Id$ is instance-independent.} By
\eqref{eq:dtheta}, $\partial\theta/\partial\F = \Id$ at $\F = \Id$ in both
columns; \S\ref{sec:fe-spaces} draws the consequence for the stability
results.

\paragraph{The instances differ in plasticity only.} The last two rows of
Table~\ref{tab:instances} are the entire difference, and they follow from
\eqref{eq:split}, as \S\ref{sec:plasticity} develops. Remark~\ref{rem:whichj2}
records that both codes implement the Simo--Hughes instance.

\subsection{Variations}
\label{sec:variations}

With $\disp \in \mathcal{U}$, the space of motions of finite deformation
energy that take the prescribed values on $\partial_\disp\Omega$, the part of
the boundary on which the motion is prescribed; $\bar\theta, \bar p \in
\mathcal{Q}$, the space of square-integrable scalar fields; test functions
$\testd \in \mathcal{U}$ vanishing on $\partial_\disp\Omega$; and
$\vartheta, \pi \in \mathcal{Q}$:
\begin{equation}
  \label{eq:variations}
  \begin{split}
    D_\disp \Pi(\testd) & =
      \int_\Omega \bar p \, \frac{\partial\theta}{\partial\disp}[\testd] \diff V
      + \int_\Omega \sdev\!\left(\edev(\disp)\right) : \frac{\partial\edev}{\partial\disp}[\testd] \diff V
      - \int_\Omega \rho_0\,~b\cdot\testd \diff V
      - \int_{\partial_T\Omega} \bT\cdot\testd \diff S = 0,
    \\
    D_{\bar\theta}\Pi(\vartheta) & =
      \int_\Omega \left( \kappa\bar\theta - \bar p \right)\vartheta \diff V = 0,
    \\
    D_{\bar p}\Pi(\pi) & =
      \int_\Omega \left( \theta(\disp) - \bar\theta \right)\pi \diff V = 0 .
  \end{split}
\end{equation}
Here $\sdev$ is the deviatoric stress conjugate to $\edev$, given by the
material law (\S\ref{sec:plasticity}). The kinematic derivatives are
\begin{equation}
  \label{eq:kinematic}
  \frac{\partial\theta}{\partial\disp}[\testd]
    = \theta'(J)\, J\, \F^{-\mathsf{T}} : \Grad\testd,
  \qquad
  \frac{\partial\edev}{\partial\disp}[\testd]
    = \mathbb{L}\supdev(\C) : \sym\!\left(\transpose{\F}\Grad\testd\right),
\end{equation}
with $\mathbb{L}(\C) := \partial \log\C/\partial\C$ the Fr\'echet derivative of
the matrix logarithm \citep{AlMohy.Higham.Relton:2013} and
$\mathbb{L}\supdev := \mathbb{L} - \tfrac13 \Id \otimes \C^{-1}$. The
volumetric derivative follows from Jacobi's formula and \eqref{eq:dtheta}. The
deviatoric expression is written for the Hencky instance of
Table~\ref{tab:instances}; the Simo--Hughes instance replaces it with the
variation of $\bar{~b}^{\mathrm e}$, and nothing outside this equation
changes.

\subsection{Stress reconstruction}
\label{sec:stress}

Collecting \eqref{eq:variations}$_1$ into the form $\int \P : \Grad\testd
\diff V$, with $\P$ the first Piola--Kirchhoff stress, the stress conjugate to
$\F$, gives in general
\begin{equation}
  \label{eq:P-general}
  \P = \theta'(J)\, J\, \bar p\, \F^{-\mathsf{T}}
     + \F \left( \mathbb{L}(\C) : \sdev \right).
\end{equation}
The deviatoric term is $\mathbb{L}$ and not $\mathbb{L}\supdev$ applied to
$\sdev$. Moving $\mathbb{L}\supdev$ off the strain variation in
\eqref{eq:kinematic} and onto $\sdev$ takes its transpose, which carries the
$\tfrac13$ term onto $\trace\sdev = 0$, and $\mathbb{L}$ is self-adjoint
because $\log\C$ is the gradient of $\trace(\C\log\C - \C)$.

When the deviatoric response is elastic, $\sdev = 2\mu\,\edev$ is a
function of the current $\C$ alone through an isotropic law, so $\sdev$ is
coaxial with $\C$ at every quadrature point; the off-diagonal blocks of
$\mathbb{L}$ in the principal basis then act on nothing and
\eqref{eq:P-general} collapses to
\begin{equation}
  \label{eq:P-simple}
  \P = \left( \varpi \Id + \sdevx \right) \F^{-\mathsf{T}} = \btau \, \F^{-\mathsf{T}},
  \qquad
  \btau := \varpi \Id + \sdevx = J\,~\sigma,
  \qquad
  \varpi := \theta'(J)\, J\, \bar p ,
\end{equation}
in which $\btau = J\,~\sigma$ is the Kirchhoff stress, $~\sigma$ the Cauchy
stress, and $\sdevx := ~R\,\sdev\,\transpose{~R}$ the spatial deviator, with
$~R$ the rotation of the polar decomposition $\F = ~R\,~U$. The material
$\sdev$ of \eqref{eq:split} and the spatial $\sdevx$ are different tensors; in
material form
the same statement reads $\P = \F^{-\mathsf T}\sdev$, which agrees with
\eqref{eq:P-simple} because $\sdev$ and $~U$ commute when coaxial.

The Kirchhoff pressure $\varpi$ carries the whole of the instance dependence.
For the Hencky instance $\theta'(J)J = 1$ and $\varpi = \bar p$; for the
Simo--Hughes instance $\theta'(J)J = J$ and $\varpi = J\bar p$, giving
$\btau = J\bar p\,\Id + \sdev$ with $\bar p = \kappa(J-1)$, which is the
Kirchhoff stress of \citet[BOX 9.1]{SimoHughes:1998} and is what both codes
compute. Agreement of \eqref{eq:P-general} with both implementations confirms
the normalization \eqref{eq:thetadef}.

\begin{remark}[Coaxiality is an elastic property]
\label{rem:coaxial}
Equation~\eqref{eq:P-simple} rests on $\sdev$ being coaxial with $\C$, which
holds only when $\sdev$ is a function of the current $\C$ alone. It fails in
two settings, and the second is the one this formulation exists for. With an
independent discrete stress field, $\sdev$ is tied to $\C$ only weakly and is
not coaxial with it at an arbitrary iterate. With plasticity, the return map
\eqref{eq:return} evaluates
$\sdev = 2\mu(\edev - \edev^{\mathrm p}_{n+1})$ against a stored
plastic strain $\edev^{\mathrm p}_n$, which is coaxial with the current $\C$
only under proportional loading; whenever the principal axes rotate ---
the large-rotation regime of \S\ref{sec:problem} --- $\sdev$ is not coaxial with $\C$, the
off-diagonal blocks of $\mathbb{L}$ do not vanish, and \eqref{eq:P-simple} is
wrong. In both settings \eqref{eq:P-simple} holds at best at a converged
proportional state, and assembling the residual from it gives a residual that
is not the derivative of the energy, with an error that is undetected under
proportional loading. The residual is assembled from \eqref{eq:P-general};
\citet{Miehe.Apel.Lambrecht:2002} carry the full projection tensor for this
reason.
\end{remark}

\begin{remark}
\label{rem:whynodev}
No deviatoric mixed pair appears because none is needed. Volumetric locking is
a constraint phenomenon: at $\kappa/\mu \gg 1$ the near-isochoric constraint
over-determines the displacement space, and relieving it requires a multiplier.
The deviatoric response is not constrained in the same sense and admits no
analogous locking mode. Adding a deviatoric pair adds degrees of freedom, an
inf-sup condition that has to be established separately, and --- if strain and
stress are placed in different spaces --- a constitutive law evaluated across
two discretizations, which is the error that corrupts plastic strain fields
when the strain is interpolated in one space and the stress evaluated in
another. \S\ref{sec:fivefield} gives the construction that adds strain as an
independent field should a later requirement demand it.
\end{remark}

\section{Discretization}

This section selects the pair of discrete spaces and gives the element-local
elimination of the auxiliary fields. $P_k$ denotes continuous piecewise
polynomials of degree $k$ on the mesh, $P_k^{\mathrm{disc}}$ the discontinuous
ones, and a bubble a polynomial on one element that vanishes on a prescribed
part of its boundary; $U_h/Q_h$ denotes the pair with displacement space
$U_h$ and pressure space $Q_h$.

\subsection{The pair}
\label{sec:fe-spaces}

Let $U_h \subset \mathcal{U}$ carry the motion and $Q_h \subset \mathcal{Q}$
carry both $\bar\theta$ and $\bar p$. Two requirements fix the choice.

\paragraph{Same space for both scalars.} $\bar\theta$ and $\bar p$ are placed in one
space $Q_h$. This is what makes the constraint term vanish identically and the
functional reduce to \eqref{eq:collapsed}; it is the only thing that reduction
needs. Given it, and a quadratic $W\supvol$ (\S\ref{sec:matreq}),
\eqref{eq:variations}$_2$ reads $\bar p = \kappa\bar\theta$ pointwise rather
than in a projected sense, and the composition
$\bar p = \kappa\,\mathcal{P}_h(\theta)$ is well defined. Splitting the two scalars
across different spaces makes the constitutive relation hold only after a
projection, determines one field only through its projection onto the other's
space, and can leave the discrete system singular.

\paragraph{Inf-sup stability, in three dimensions.} By \eqref{eq:kinematic} the
pressure--displacement pairing is
$\int \varpi \,\F^{-\mathsf T}\!:\!\Grad\testd$ with
$\varpi = \theta'(J)\,J\,\bar p$, and not $\int \bar p
\,\F^{-\mathsf T}\!:\!\Grad\testd$; the two coincide at $\F = \Id$ and nowhere
else. The normalization \eqref{eq:thetadef} sets $\theta'(1) = 1$, so
$\varpi = \bar p$ there and the pairing reduces to the divergence pairing
$\int \bar p \Div\testd \diff V$ --- in both instances of
Table~\ref{tab:instances}, which is what makes the stability results of
\S\ref{sec:beta} a property of the pair rather than of the material. That
pairing is the one of incompressible linear elasticity, in which the pressure
is the multiplier of the constraint $\Div\disp = 0$, and its discrete
inf-sup condition,
\begin{equation}
  \label{eq:infsup}
  \inf_{q \in Q_h}\;\sup_{~v \in U_h}\;
  \frac{\int_\Omega q \, \Div ~v \diff V}{\lvert ~v \rvert_{H^1} \norm{q}_{L^2}}
  \;\ge\; \beta_h ,
\end{equation}
is the condition on the pair. The admissible pairs are therefore those for
which $\beta_h$ in \eqref{eq:infsup} is bounded below uniformly in $h$, and
stability must also hold on the target range of $\kappa/\mu$
\citep{BoffiBrezziFortin:2013}.

Stability of a pair depends on the spatial dimension. The conforming Crouzeix--Raviart pair \citep{CrouzeixRaviart:1973} is, in two
dimensions, $P_2$ enriched with a cubic interior bubble over
$P_1^{\mathrm{disc}}$; its three-dimensional counterpart below carries a
quartic interior bubble and cubic face bubbles, and the quadrature degree the
element needs follows from the quartic. Unenriched
$P_2/P_1^{\mathrm{disc}}$ is stable in two dimensions only on special meshes.
In three dimensions \citet[\S8.7]{BoffiBrezziFortin:2013} state that a
constant discontinuous pressure is not controlled by $P_2$: ``to control
piecewise constant pressure, we need some degrees of freedom on the faces,''
and the quadratic space has none. The stable three-dimensional choices with a
discontinuous pressure therefore require bubble enrichment of the displacement
space, and the enrichment has two parts with two functions: face bubbles control the
constant pressure, and an interior bubble lifts it to a linear one
\citep[Example~8.7.2]{BoffiBrezziFortin:2013}. Table~\ref{tab:pairs} lists the
candidates, and \S\ref{sec:beta} measures every one of them.

\begin{table}[h]
\centering
\begin{tabular}{@{}llll@{}}
\toprule
Displacement & Pressure & 3D stable & Pressure elimination \\
\midrule
$P_1$                        & $P_0$                   & no  & local \\
$P_2$                        & $P_0$                   & \textbf{no} & local \\
$P_2 \oplus b_{\mathrm{face}}$ & $P_0$                 & yes & local \\
$P_2$                        & $P_1^{\mathrm{disc}}$   & \textbf{no} & local \\
$P_2 \oplus b_{\mathrm{int}}$ & $P_1^{\mathrm{disc}}$  & no  & local \\
$P_2 \oplus b_{\mathrm{int}} \oplus b_{\mathrm{face}}$ & $P_1^{\mathrm{disc}}$ & yes & local \\
$P_2$                        & $P_1$ continuous        & yes & global \\
\bottomrule
\end{tabular}
\caption{Displacement--pressure pairs on tetrahedra, with three-dimensional
stability as \citet[\S8.7, Example~8.7.2]{BoffiBrezziFortin:2013} give it.
Local elimination
requires a discontinuous pressure; three-dimensional stability with a
discontinuous pressure requires face bubbles, and a linear pressure requires
an interior bubble as well.}
\label{tab:pairs}
\end{table}

The pair this formulation requires is $P_2$ enriched with face and
interior bubbles over $P_1^{\mathrm{disc}}$. The two enrichments are different
kinds of object. The interior bubble vanishes on the whole element boundary: it is never shared,
never reached by a boundary condition, and condenses away, adding no global
unknown. A face bubble vanishes on three faces and not on its own, so the two
elements meeting there share it; it is a global unknown and does not condense.
The enrichment therefore adds global degrees of freedom, which
\S\ref{sec:infsup} counts. The pressure elimination of \S\ref{sec:local}
nevertheless stays element by element, because it rests on $Q_h$ being
discontinuous and not on the displacement space.

The linear pressure is not forced by the stability results. $P_2 \oplus
b_{\mathrm{face}}/P_0$ is also stable (\S\ref{sec:beta}), at a higher constant
and with the same number of global unknowns, since the interior bubble condenses. The count
of \S\ref{sec:firstresult} applies to the unenriched displacement space, and
\S\ref{sec:assembled} shows that its modes do not persist after assembly. What
separates the two pressure spaces is the claim of
\S\ref{sec:prediction}: a constant pressure leaves in $\ker ~G_h$ the
element-scale dilatational modes that a plastic tangent can free, and a linear
one does not. \S\ref{sec:plastic} measures both pairs, and that measurement
is the reason for the linear pressure.

Both that pair and the Taylor--Hood pair --- continuous $P_2$ over continuous
$P_1$, whose stability in three dimensions is established
\citep{BoffiBrezziFortin:2013} and which serves as the positive control of
the measurements --- are measured in \S\ref{sec:beta}, and every other
element-local pair decays. The difference between the two stable pairs is not
stability but where the coupling is --- a global pressure that cannot be
eliminated element by element against global face bubbles that cannot be
condensed --- and \S\ref{sec:risks} quantifies it.

\subsection{Local elimination}
\label{sec:local}

Because $Q_h$ is discontinuous, the mass matrix
$[M]_{mn} = \int_\Omega \chi_m \chi_n \diff V$, with $\{\chi_m\}$ a basis of
$Q_h$, is block diagonal by element and constant in the reference
configuration. One factorization at setup serves
every Newton iteration of every step, and
\begin{equation}
  \label{eq:projection}
  M \{\bar\theta\} = \int_\Omega \chi_m \, \theta(\disp) \diff V,
  \qquad
  \bar p = \kappa \bar\theta ,
\end{equation}
is an element-local solve, with $\theta = \log J$ for a Hencky material and
$\theta = J - 1$ for the Simo--Hughes one (Table~\ref{tab:instances}). The second
relation in \eqref{eq:projection} holds pointwise, rather than only
after projection, precisely because $W\supvol$ is quadratic; for a material
where it is not, it is replaced by a second projection with the same
factorization (Remark~\ref{rem:nonquadratic}). Substituting $\bar p(\disp)$ into
\eqref{eq:variations}$_1$ leaves a displacement-sized nonlinear problem whose
tangent retains the sparsity of the displacement mesh --- the auxiliary fields
never couple two elements. This is the property that keeps the composite
tetrahedron's tangent as sparse as a displacement element's, and it is
preserved here.

\begin{remark}
The locality is a consequence of the discontinuous pressure, not of the
Hu--Washizu structure. A conforming auxiliary field --- one requiring
inter-element continuity --- would make $\bar p(\disp)$ a global function of
$\disp$, and every tangent--vector product would then require a global solve
with its mass matrix, which the composite tetrahedron does not require. That
is the main reason the pressure is kept discontinuous.
\end{remark}

\section{Plasticity}
\label{sec:plasticity}

This section develops the logarithmic instance of Table~\ref{tab:instances};
the Simo--Hughes column differs here and nowhere else (\S\ref{sec:instances}),
and Remark~\ref{rem:whichj2} states which model the open measurement of
\S\ref{sec:verify} uses.

The deviatoric response is evaluated at quadrature points of the displacement
element, on $\edev(\disp)$ computed there. The plastic strain is carried as a
logarithmic strain $\edev^{\mathrm p}$ and subtracted additively --- the
construction of \citet{Miehe.Apel.Lambrecht:2002}, in which a material
geometric pre-processing $\C \mapsto \eps$ and a post-processing through
$\mathbb{L}(\C)$ precede and follow an unmodified small-strain algorithm. This is a
different constitutive model from the multiplicative $\F = \F^{\mathrm e}
\F^{\mathrm p}$ plasticity of \citet{Simo:1992}, which also preserves the
small-strain return map but applies it to the elastic logarithmic strain of a
trial elastic state with an exponential update; the two coincide under
proportional loading with fixed principal axes and differ when the axes rotate.
Because
the constitutive law is linear in logarithmic strain, the $J_2$ return map is
the small-strain radial return --- the return map that scales the trial
deviatoric stress toward the yield surface along its own direction ---
verbatim:
\begin{equation}
  \label{eq:return}
  \begin{split}
    \edev^{\mathrm{tr}} &:= \edev(\disp) - \edev^{\mathrm p}_n,
    \qquad
    \sdev^{\mathrm{tr}} := 2\mu\,\edev^{\mathrm{tr}},
    \\
    \Delta\gamma &: \quad
      \norm{\sdev^{\mathrm{tr}}} - 2\mu\,\Delta\gamma
      - \sqrt{2/3}\left[\sigma_y + \mathcal{H}'\!\left(\bar\epsilon^{\mathrm p}_n + \sqrt{2/3}\,\Delta\gamma\right)\right] = 0
      \quad\text{if } f^{\mathrm{tr}} > 0, \quad \text{else } \Delta\gamma := 0,
    \\
    \edev^{\mathrm p}_{n+1} &:= \edev^{\mathrm p}_n + \Delta\gamma\, ~n,
    \qquad ~n := \sdev^{\mathrm{tr}} / \norm{\sdev^{\mathrm{tr}}},
    \qquad
    \bar\epsilon^{\mathrm p}_{n+1} := \bar\epsilon^{\mathrm p}_n + \sqrt{2/3}\,\Delta\gamma,
    \\
    \sdev &:= 2\mu\left(\edev(\disp) - \edev^{\mathrm p}_{n+1}\right).
  \end{split}
\end{equation}
Here $\sigma_y$ is the initial yield stress, $\bar\epsilon^{\mathrm p}$ the
equivalent plastic strain, $\mathcal{H}'$ the hardening function (the increase
of the yield stress with $\bar\epsilon^{\mathrm p}$) and $\mathcal{H}''$ its
derivative, the hardening modulus, and
$f^{\mathrm{tr}} := \norm{\sdev^{\mathrm{tr}}} - \sqrt{2/3}\,[\sigma_y +
\mathcal{H}'(\bar\epsilon^{\mathrm p}_n)]$ the trial yield function. The
consistency condition is a scalar equation in $\Delta\gamma$. For linear
hardening it closes in one step,
$\Delta\gamma = f^{\mathrm{tr}}/(2\mu + \tfrac23\mathcal{H}'')$, which is the
form both codes implement; for nonlinear hardening it is a scalar Newton
iteration \citep[BOX~3.2]{SimoHughes:1998}. Internal variables are stored at quadrature points and
are never projected, so no interpolation error enters the history. Frame-indifferent interpolation, when
tensorial internal variables must be transferred, follows
\citet{Mota.etal:2013}.

Three properties follow. They are the motivation for the formulation rather
than incidental consequences of it.

\paragraph{The constitutive law is untouched.} All finite-deformation content is
absorbed into two material-independent geometric transforms: the forward
map $\C \mapsto \eps$ and the tangent map $\mathbb{L}(\C)$
\citep{Miehe.Apel.Lambrecht:2002}. Between them sits an unmodified small-strain
algorithm. Any model already written for small strain --- $J_2$, kinematic
hardening, viscoplasticity, damage --- is usable without reformulation. This is the
practical content of preferring a Hu--Washizu structure over one built around a
particular stress--strain pairing. This property belongs to the
logarithmic measure and not to the Hu--Washizu structure, so it is the one
place where \S\ref{sec:matreq} does not apply: an admissible material
in the sense of \eqref{eq:matsplit} need not have it, and the Simo--Hughes
model both codes implement does not (Remark~\ref{rem:whichj2}).

\paragraph{Strain and stress never cross discretizations.} $\edev$ is evaluated and
$\sdev$ returned at the same quadrature point of the same element. There is no
projection between the two, so the interpolation error between two discrete
spaces that corrupts the equivalent plastic strain in stabilized formulations
has no mechanism here.

\paragraph{No stabilization enters the deviator.} There is no subgrid term, no
penalty and no tuned coefficient anywhere in the formulation, so nothing of the
form \eqref{eq:vms} can reach the return map.

\begin{remark}[The model implemented in both codes, and the open measurement]
\label{rem:whichj2}
Equation \eqref{eq:return} is the return map this formulation specifies.
Neither Norma nor Carina implements it. Both implement the Simo--Hughes $J_2$ model
\citep{SimoHughes:1998}: a neo-Hookean volumetric--isochoric split,
$W = \tfrac{\kappa}{2}(J-1)^2 + \tfrac{\mu}{2}(\trace\bar{~b}^{\mathrm e} - 3)$,
with the radial return taken on the isochoric elastic left Cauchy--Green tensor
rather than on logarithmic strain. It satisfies \eqref{eq:matsplit} exactly ---
$\det\bar{~b}^{\mathrm e} = 1$ makes the isochoric term independent of $J$ by
construction rather than by an identity --- and its $W\supvol$ is quadratic in
$J-1$, so it is admissible here in the variable $J-1$.

Two consequences follow. The first is that the smoothness of the equivalent
plastic strain in a nonlinear solve, the open measurement of
\S\ref{sec:verify}, is measured with Simo--Hughes unless a logarithmic $J_2$
is written for the purpose, and the substitution is larger than an elastic
split: it exchanges the additive logarithmic model of \eqref{eq:return} for
multiplicative plasticity with the return map on $\bar{~b}^{\mathrm e}$. The
smoothness of equivalent plastic strain is a question about whether strain and
stress cross discretizations, which is independent of the plasticity model
carrying them; the substitution changes the plasticity model and not only the
volumetric split.

The second is that the argument above that small-strain models are usable
unmodified is weakened by it. That property belongs to the
logarithmic measure, and Simo--Hughes does not have it: its return map is
specific to its own structure. The argument stands for the formulation as
specified and does not describe what is implemented.
\end{remark}

\section{The claim}
\label{sec:prediction}

This section states the claim of the note: which hypothesis of Brezzi's
theorem a plastic tangent can violate, and what distinguishes pairs under it.
The second hypothesis is coercivity of the deviatoric form on $\ker ~G_h$
(\S\ref{sec:problem}).

In elasticity the second hypothesis holds for every conforming pair: by
Korn's inequality, which bounds $\int \lvert\Grad ~v\rvert^2 \diff V$ by a
multiple of $\int \lvert\dev\eps(~v)\rvert^2 \diff V$ for $~v \in H^1_0$,
the deviatoric form is coercive on all of $U_h$, so it is coercive on
$\ker ~G_h$, the set of discrete displacements whose divergence has zero
projection onto $Q_h$ (\S\ref{sec:method}), and no elastic soft mode is a
failure of the theorem.
A count of exact zero-energy modes on one element is therefore not an argument
against a constant pressure in elasticity; \S\ref{sec:assembled} shows that
those modes do not persist after assembly.

Under plastic loading the consistent $J_2$ tangent, the derivative of the
stress returned by \eqref{eq:return} with respect to the strain, is
\begin{equation}
  \label{eq:cep}
  \mathbb{C}^{\mathrm{ep}} = 2\mu \left( \mathbb{I}\supdev - \eta\, ~n \otimes ~n \right),
  \qquad
  \mathbb{I}\supdev := \mathbb{I} - \tfrac13 \Id \otimes \Id,
  \qquad
  ~n := \sdev^{\mathrm{tr}} / \norm{\sdev^{\mathrm{tr}}},
  \qquad
  \eta = \frac{2\mu}{2\mu + \tfrac{2}{3}\mathcal{H}''} ,
\end{equation}
with $\eta = 1$ for perfect plasticity: the shear stiffness along the flow
direction is removed.

Read as a constant, the second hypothesis cannot separate pairs. For every $~v \in H^1_0(\Omega)$,
integration by parts gives $\int \lvert\dev\eps\rvert^2 \diff V =
\tfrac12 \int \lvert\Grad ~v\rvert^2 \diff V + \tfrac16 \int (\Div ~v)^2
\diff V$, and $(~n : \eps)^2 \le \lvert\dev\eps\rvert^2$ since $~n$ is a
unit deviator, so
\begin{equation}
  \label{eq:floor}
  \frac{1}{2\mu} \int \eps : \mathbb{C}^{\mathrm{ep}} : \eps \diff V
  \;\ge\; (1 - \eta) \int \lvert\dev\eps\rvert^2 \diff V
  \;\ge\; \frac{1 - \eta}{2} \int \lvert\Grad ~v\rvert^2 \diff V
\end{equation}
for every conforming pair and every member of its kernel, isochoric or not.
For $\eta < 1$ the coercivity constant on $\ker ~G_h$ is therefore bounded
below uniformly in $h$ and in the pair. For $\eta = 1$ the continuum
operator loses strong ellipticity, positivity of the energy of every
displacement jump across a plane: a shear band, a displacement jump
$~e_1$ across a plane with normal $~e_2$, has $\eps \propto \sym(~e_1 \otimes ~e_2)$, to which \eqref{eq:cep} with
that $~n$ assigns zero energy, so the constant tends to zero under refinement for
every pair as the discrete space resolves narrower bands
(\S\ref{sec:plastic}). Neither limit says anything about the pair.

What distinguishes the pairs is what the soft directions are. A member of
$\ker ~G_h$ that is nearly isochoric is a shear band the continuum problem
also has, and its low energy is physical. A member with pointwise nonzero
dilatation of zero element mean has the volumetric energy
$\kappa \int (\Div ~v)^2 \diff V$ in the continuum and none in the element,
so once the tangent removes its shear energy the element assigns it zero
energy. That is the soft
mode reported for constant-pressure elements, and whether $\ker ~G_h$
contains such members is a property of the pair. For a pair satisfying the
inf-sup condition, $\ker ~G_h$ approximates the space of isochoric
displacements at the rate of $U_h$ \citep[Prop.~5.1.3]{BoffiBrezziFortin:2013},
so the dilatation carried by any of its members vanishes with $h$. For an
element-constant pressure on a quadratic displacement, $\ker ~G_h$ contains
displacements whose dilatation is nonzero pointwise and averages to zero over
each element, and their dilatation does not vanish with $h$.

The claim is therefore the following. \emph{Under a plastic tangent, every
soft mode of the stable pair with a linear discontinuous pressure carries a
dilatation $\int (\Div ~v)^2 \diff V / \int \lvert\Grad ~v\rvert^2 \diff V$
that vanishes under refinement, and a constant pressure has soft modes whose
dilatation does not; the soft modes reported for constant-pressure elements
are the latter, and no stabilization is needed where they are absent.} Both operators are linear at a fixed plastic state,
so the claim is tested without a nonlinear solve, and \S\ref{sec:plastic}
tests it. If the dilatation of the soft modes does not separate the pairs,
the formulation has no advantage over the constant-pressure elements it is
meant to replace.

\section{Dynamics}
\label{sec:dynamics}

The spatial formulation is unchanged; the mass matrix enters only the
displacement block. Two implementation choices are specified.

\paragraph{Lumping.} Row-sum lumping, which assigns to each degree of freedom
the row sum of the consistent mass matrix, produces non-positive vertex masses
on the quadratic tetrahedron and cannot be used. Diagonal-scaling (HRZ) lumping
\citep{HintonRockZienkiewicz:1976}, which scales the diagonal of the consistent
mass to preserve the total mass, is positive for any basis, so positivity is
not the limitation; on quadratic simplices the known disadvantage of HRZ is
accuracy, with the vertex masses small. The limitation the enrichment
introduces depends on the basis that carries it, and \S\ref{sec:basis}
measures both choices. In the hierarchical basis --- the ten Lagrange
functions of the quadratic tetrahedron together with the five bubbles --- a
rigid translation has zero coefficient on every bubble degree of freedom, so
their diagonal mass is inert for translations, yet HRZ's normalization counts
it: the Lagrange masses sum to less than the total mass by the fraction
assigned to the bubbles, $51\%$ on the reference tetrahedron, and a rigid
translation carries half its momentum. The normalization must then run
over the Lagrange degrees of freedom alone, with each bubble keeping its own
consistent diagonal, which restores the momentum exactly. With a nodal basis
on the fifteen nodes of the \code{TETRA15} arrangement --- vertices, edge
midpoints, face centroids, centroid --- the functions form a partition of
unity, HRZ conserves momentum without qualification, and the lumped masses
are positive on every element measured. Row-sum lumping was also positive in
that basis on the three elements measured, unlike for the ten-node quadratic
tetrahedron (\code{TETRA10}), but its positivity is not guaranteed, and HRZ
lumping is used.

\paragraph{Critical time step.} From the element eigenvalue bound at
$\kappa/\mu = 2$ (\S\ref{sec:basis}), the Courant number $c_p\,\Delta t/h$,
with $c_p = \sqrt{(\kappa + 4\mu/3)/\rho}$ the dilatational wave speed,
$\rho$ the density, $\Delta t$ the stable step and $h$ the element edge
length, is given in Table~\ref{tab:basis} for the four-node linear
tetrahedron (\code{TETRA4}), for \code{TETRA10} under HRZ, and for the
enriched element in both bases. The quadratic element reduces the stable step
by a factor of $2.4$ against the linear one on the same mesh, and the
enrichment a further $1.4$ with the nodal basis; the ordering is the same on
the distorted element and on a Kuhn tetrahedron of the Freudenthal mesh
(\S\ref{sec:method}). A $P_2$ mesh resolves a given feature with fewer
elements; whether the smaller element count compensates for the smaller step
depends on the problem and is to be measured.

\section{Optional five-field variant}
\label{sec:fivefield}

This section gives the extension to an independent deviatoric strain field,
for non-local regularization or for internal variables requiring their own
discrete representation: it introduces
$\bar\edev$ and $\bar\sdev$ in the same discontinuous space, of the same
degree as $Q_h$. Sharing the space keeps the return map at shared
quadrature points, keeps every elimination element-local, leaves the discrete
system non-singular, and --- for a linear deviatoric law, and only then ---
makes the constitutive relation hold pointwise rather than after a projection;
with plasticity $\bar\sdev = \mathcal{P}_h\,\sdev(\bar\edev)$ is projected. Placing $\bar\edev$ and $\bar\sdev$ in different spaces loses
each of these properties. Remark~\ref{rem:coaxial} also applies: with an
independent $\bar\sdev$, the residual must be assembled from
\eqref{eq:P-general} and not from \eqref{eq:P-simple}.

\section{Numerical results}
\label{sec:results}

Every measurement below is on a linear operator, and all but those of
\S\ref{sec:deformed} are at $\F = \Id$: the condensed operator
\eqref{eq:condensed} for the count and the locking sweep, the $H^1$ seminorm
for the stability sweep, and the deviatoric form under a plastic tangent for
the soft-mode measurements. None needs a solver, a material model,
or a mesh reader. The scripts are in \code{prototype/}, each reports its own
result, and \code{checks.jl} verifies the quadrature, the bubbles, and the
conformity of the enriched spaces against values known in advance.

\subsection{Method}
\label{sec:method}

Meshes are generated rather than read, so the refinement sequence is exact.
The unit cube is divided into $N^3$ cubes and each cube into six tetrahedra by
the Freudenthal (Kuhn) subdivision --- the six permutations $(a,b,c)$ giving
the Kuhn tetrahedra $\{~0, ~e_a, ~e_a + ~e_b, (1,1,1)\}$. That decomposition
is conforming across shared cube faces, which the five-tetrahedron split is
not. Quadratic meshes add edge midpoints in the node
order the reference element expects.

For the stability test the displacement is fixed on the entire boundary and
three operators are assembled,
\begin{equation}
  \label{eq:infsup-ops}
  K_{ij} = \int \Grad N_i : \Grad N_j \diff V,
  \qquad
  G_{mi} = \int \chi_m \, \Div N_i \diff V,
  \qquad
  M_{mn} = \int \chi_m \chi_n \diff V,
\end{equation}
with $N_i$ the displacement basis functions, $\chi_m$ the pressure basis
functions, and $K$ the $H^1$ seminorm, the norm on the displacement in
\eqref{eq:infsup}; $\ker ~G$ is the discretely isochoric subspace, the set of
displacements whose divergence has zero projection onto $Q_h$. The discrete inf-sup
constant is $\beta_h = \sqrt{\lambda_{\min}}$ over the nonzero spectrum of
\begin{equation}
  \label{eq:infsup-eig}
  ~G ~K^{-1} \transpose{~G}\, q = \lambda ~M q .
\end{equation}
Boundary data fixed everywhere always leaves one zero mode, the hydrostatic
one, because $\int \Div ~v = 0$ for $~v \in H^1_0$; any zero beyond
that is a spurious pressure mode.

\subsection{Spurious modes on one element}
\label{sec:firstresult}

For linear response the count of zero-energy modes is a rank count on a single
element. Condensing the three fields gives
\begin{equation}
  \label{eq:condensed}
  ~K_e = ~K\supdev + ~K\supvol,
  \qquad
  ~K\supvol := \kappa\, \transpose{~G} ~M^{-1} ~G,
  \qquad
  M_{mn} = \int_e \chi_m \chi_n \diff V,
  \qquad
  G_{mi} = \int_e \chi_m \left(\trace ~B\right)_i \diff V ,
\end{equation}
with $~B$ the strain--displacement operator of the element and $m = \dim Q_h$
on one element.
$~K\supdev$ annihilates every mode with zero deviatoric strain at quadrature:
the six rigid-body modes, plus the volumetric directions the displacement space
can produce. The second term has rank equal to the number of those directions
the pressure space resists, at most $m$. Anything left over carries neither
deviatoric nor volumetric energy and is exactly zero-energy:
\begin{equation}
  \label{eq:spurious}
  n_{\mathrm{spurious}} = \dim \ker ~K\supdev - 6 - \operatorname{rank} ~K\supvol .
\end{equation}

For the quadratic tetrahedron $\dim\ker ~K\supdev = 10$: six rigid modes and
four volumetric directions, the latter because $\trace\eps(\disp)$ is a $P_1$
field. Table~\ref{tab:count} reports \eqref{eq:spurious} against the measured
spectrum.

\begin{table}[h]
\centering
\begin{tabular}{@{}lccccc@{}}
\toprule
Pressure space & $m$ & $\dim\ker ~K\supdev$ & $\operatorname{rank} ~K\supvol$ & predicted & observed \\
\midrule
$P_0$ constant            & 1  & 10 & 1 & \textbf{3} & \textbf{3} \\
$P_1$ discontinuous       & 4  & 10 & 4 & 0 & 0 \\
$P_2$ discontinuous       & 10 & 10 & 4 & 0 & 0 \\
\bottomrule
\end{tabular}
\caption{Spurious zero-energy modes of the three-field quadratic tetrahedron at
$\kappa/\mu = 10^4$, on the affine element. Script:
\code{prototype/softmode.jl}.}
\label{tab:count}
\end{table}

Four consequences follow.

The constant-pressure element carries three modes that are exactly
zero-energy on one element, not merely soft. Whether they persist after assembly is
the question of \S\ref{sec:assembled}.

A $P_1$-discontinuous pressure removes them without an added energy term and
leaves the six rigid-body modes. That is an element-level fact:
\S\ref{sec:assembled} and \S\ref{sec:beta} show that the same pair locks and
is unstable once assembled.

Enriching the pressure further changes nothing: the rank saturates at four,
because four is all the displacement space can produce. Enriching the
pressure beyond $P_1^{\mathrm{disc}}$ adds unknowns and leaves the rank
unchanged.

The count is a property of the affine element. On a curved element, with the
midside nodes displaced so that $\det ~J$ ranges over $[0.55, 1.19]$, the
isoparametric displacement space is not polynomial in $\bx$, only the
dilatation of its four volumetric directions remains free of deviatoric
strain, $\dim\ker ~K\supdev = 7$, and the three constant-pressure modes are
soft rather than exactly zero: the softest nonzero mode of the
constant-pressure element lies one to two orders of magnitude below the
softest mode of the linear-pressure element on the same curved element. A
four-point rule samples a curved element at four points and reports the
affine counts there as well; the count uses a $27$-point rule. The assembled
behavior of \S\ref{sec:assembled} rests on soft rather than exact modes.

\begin{remark}
This identifies the mechanism. Inf-sup stability is a
global property of a mesh sequence and is invisible on one element; the rank
argument above says which modes exist, not whether the constant in the stability
estimate persists under refinement; \S\ref{sec:beta} measures it.
\end{remark}

\subsection{A necessary condition from degree-of-freedom counts}
\label{sec:infsup}

Whether the constant in the stability estimate persists under refinement is a
property of a mesh sequence, and a necessary condition for it follows from
counting.

Write $n_u$ for the number of free displacement degrees of freedom and $n_p$
for the pressure degrees of freedom. The Schur complement of the mixed system
is $~S = ~G ~K^{-1} \transpose{~G}$, an $n_p \times n_p$ matrix of rank at most
$n_u$. If $n_p > n_u$ then $~S$ is singular by counting, the pressure
space carries modes no displacement can resist, and no inf-sup constant exists
--- regardless of the values those constants would otherwise take.

Table~\ref{tab:dofcount} counts them on a Freudenthal-subdivided cube.

\begin{table}[h]
\centering
\begin{tabular}{@{}lrrrr@{}}
\toprule
Pair & $N$ & $n_u$ & $n_p$ & $n_p/n_u$ \\
\midrule
$P_2/P_0$                              & 12 & 36{,}501 & 10{,}368 & 0.284 \\
$P_2/P_1^{\mathrm{disc}}$              & 12 & 36{,}501 & 41{,}472 & \textbf{1.136} \\
$P_2 \oplus b_{\mathrm{int}} / P_1^{\mathrm{disc}}$ & 12 & 67{,}605 & 41{,}472 & 0.613 \\
$P_2 \oplus b_{\mathrm{int}} \oplus b_{\mathrm{face}} / P_1^{\mathrm{disc}}$
                                       & 12 & 127{,}221 & 41{,}472 & 0.326 \\
\bottomrule
\end{tabular}
\caption{Degree-of-freedom counts on the unit cube, displacement fixed on the
boundary. Script: \code{prototype/infsup.jl}.}
\label{tab:dofcount}
\end{table}

For $P_2/P_1^{\mathrm{disc}}$ the ratio exceeds one, and its limit is
computable. The quadratic nodes of this mesh family fill the $(2N+1)^3$
grid, so with the whole boundary fixed $n_u = 3(2N-1)^3 \to 24N^3$, while
$n_p = 4 \cdot 6N^3 = 24N^3$ regardless of boundary data. The ratio tends to
one: the pair is critically determined in the limit, which excludes a uniform
inf-sup constant.

The finite-$N$ ratio depends on the boundary condition: with the whole
boundary fixed, $n_u = 3(2N-1)^3$ and the ratio exceeds one at every $N$; with
one face fixed, $n_u = 6N(2N+1)^2$ and it is below one at every $N$. Both tend
to one.

\begin{proposition}
On a tetrahedral mesh family with at least six tetrahedra per vertex
asymptotically --- the Freudenthal family among them --- and the boundary
fixed, $P_2/P_1^{\mathrm{disc}}$ is not inf-sup stable, and no refinement
rescues it.
\end{proposition}

The restriction to the mesh family is necessary. On a general tetrahedral mesh
with $n_V$ vertices, $n_E$ edges and $n_T$ tetrahedra, Euler's relation for a
simply connected three-dimensional mesh gives $n_E \approx n_V + n_T$, so
$n_u \approx 3(2n_V + n_T)$, $n_p = 4n_T$, and
$n_p/n_u \to 4/(3(2n_V/n_T + 1))$, which exceeds one only when
$n_T/n_V > 6$. Below that the count refutes nothing, and on other mesh families
the pair can be stable: \citet{Zhang:2011} shows $P_2/P_1^{\mathrm{disc}}$ stable on Powell--Sabin
tetrahedral meshes once their known spurious pressure modes are filtered. The
general three-dimensional result is \citet[Remark~8.6.2]{BoffiBrezziFortin:2013}:
with a $P_k$ displacement, a discontinuous pressure is stable on general meshes
only at degree $k - 3$, so for $k = 2$ nothing is admissible without enrichment.

Adding one interior bubble per element supplies $3n_{\mathrm{el}}$ further
displacement degrees of freedom and moves the ratio to $0.613$; adding the four
face bubbles as well --- the enrichment \S\ref{sec:fe-spaces} calls for ---
supplies $3n_{\mathrm{f}}$ more, one vector unknown per interior face,
and moves it to $0.326$. Both satisfy $n_p < n_u$. The enrichment also fixes
the number of global unknowns. The interior bubble
condenses; the face bubbles are shared and do not. The global displacement
count of the full enrichment is therefore $36{,}501 + 59{,}616 = 96{,}117$,
which is $2.6$ times the $36{,}501$ of $P_2/P_0$ --- also the nodal count of
the composite tetrahedron on the same \code{TETRA10} mesh --- and $2.5$ times
the $38{,}698$ global unknowns of Taylor--Hood, whose $2{,}197$ vertex pressures
cannot be eliminated. The $127{,}221$ in Table~\ref{tab:dofcount} is the
dimension of the space, not the number of global unknowns.

The count has precedents. $n_p/n_u$ is the constraint count of
\citet{NagtegaalParksRice:1974}, whose subject --- locking of simplicial
elements in the fully plastic range --- is this note's, and its reciprocal
$\varrho = n_u/n_p$ is the constraint ratio of \citet[\S4.3.7]{Hughes:1987},
with $\varrho = 3$ the optimum in three dimensions. At $N = 12$:
$\varrho = 3.52$ for $P_2/P_0$, $0.88$ for $P_2/P_1^{\mathrm{disc}}$, $1.63$
with the interior bubble, and $3.07$ for the full enrichment. The stable pair
has the optimal value, and the isochoric fraction $1 - n_p/n_u$ of
\S\ref{sec:assembled} is $1 - 1/\varrho$, with $2/3$ the optimum.

Whether that limiting degeneracy is the decisive objection on the meshes
tested is a separate question, which \S\ref{sec:assembled} answers.

\subsection{Assembled behavior}
\label{sec:assembled}

The element count of \S\ref{sec:firstresult} is a local statement. An
element-local mode persists after assembly only if its trace is compatible
with every neighbor across every shared face, so the global count can be
smaller than $3n_{\mathrm{el}}$, and it is what the argument against a
constant pressure rests on.

Table~\ref{tab:assembled} assembles $~K\supdev + \kappa \transpose{~G} ~M^{-1}
~G$ over meshes under two boundary conditions: the whole boundary fixed, and
the base $z = 0$ alone fixed. The second is closer to the conditions of a body confined
on part of its boundary and free elsewhere, under which soft modes are
reported.

The locking measure is again a rank. A displacement mode has volumetric
energy exactly when it lies outside $\ker ~G$, so as $\kappa \to \infty$ the
reachable deformations collapse onto that kernel and
$\dim\ker ~G = n_u - \operatorname{rank} ~G$ is the discretely isochoric
subspace. A large fraction means the element still deforms; a collapse toward
zero is volumetric locking.

The fraction is compared with a bound. Since
$\operatorname{rank}~G \le n_p$, it cannot fall below $1 - n_p/n_u$. A pair
that attains that bound has $\operatorname{rank}~G = n_p$: every pressure mode
is resisted by a distinct displacement direction. A fraction above the bound
means that $n_p - \operatorname{rank}~G$ pressure modes are unresisted; a pair
that cannot reach it, because $n_p > n_u$, is over-determined. The difference $n_p - \operatorname{rank}~G$ counts the
pressure modes no displacement resists, which with the whole boundary fixed
should be exactly one, the hydrostatic mode; Table~\ref{tab:beta} carries it
as a column.

\begin{table}[h]
\centering
\begin{tabular}{@{}llrrrrrr@{}}
\toprule
boundary & Pair & $N$ & $n_u$ & $n_p$ & $\operatorname{rank}~G$ & fraction & bound \\
\midrule
whole & $P_2/P_0$                 & 3 & 375  & 162  & 158  & 0.579 & 0.568 \\
      &                           & 4 & 1029 & 384  & 380  & 0.631 & 0.627 \\
      & $P_2/P_1^{\mathrm{disc}}$ & 3 & 375  & 648  & 363  & \textbf{0.032} & --- \\
      &                           & 4 & 1029 & 1536 & 957  & \textbf{0.070} & --- \\
      & $\oplus\, b_{\mathrm{int}}$
                                  & 3 & 861  & 648  & 644  & 0.252 & 0.247 \\
      &                           & 4 & 2181 & 1536 & 1532 & 0.298 & 0.296 \\
      & $\oplus\, b_{\mathrm{face}}$
                                  & 3 & 1185 & 648  & 646  & 0.455 & 0.453 \\
      &                           & 4 & 3045 & 1536 & 1534 & 0.496 & 0.496 \\
      & $\oplus\, b_{\mathrm{int}}, b_{\mathrm{face}}$
                                  & 3 & 1671 & 648  & 647  & \textbf{0.613} & 0.612 \\
      &                           & 4 & 4197 & 1536 & 1535 & \textbf{0.634} & 0.634 \\
      & $P_2 \oplus b_{\mathrm{face}}/P_0$
                                  & 3 & 1185 & 162  & 161  & 0.864 & 0.863 \\
      &                           & 4 & 3045 & 384  & 383  & 0.874 & 0.874 \\
\midrule
base  & $P_2/P_0$                 & 3 & 882  & 162  & 162  & 0.816 & 0.816 \\
      &                           & 4 & 1944 & 384  & 384  & 0.802 & 0.802 \\
      & $P_2/P_1^{\mathrm{disc}}$ & 3 & 882  & 648  & 522  & 0.408 & 0.265 \\
      &                           & 4 & 1944 & 1536 & 1216 & 0.374 & 0.210 \\
      & $\oplus\, b_{\mathrm{int}}$
                                  & 3 & 1368 & 648  & 648  & 0.526 & 0.526 \\
      &                           & 4 & 3096 & 1536 & 1536 & 0.504 & 0.504 \\
      & $\oplus\, b_{\mathrm{face}}$
                                  & 3 & 1962 & 648  & 648  & 0.670 & 0.670 \\
      &                           & 4 & 4440 & 1536 & 1536 & 0.654 & 0.654 \\
      & $\oplus\, b_{\mathrm{int}}, b_{\mathrm{face}}$
                                  & 3 & 2448 & 648  & 648  & 0.735 & 0.735 \\
      &                           & 4 & 5592 & 1536 & 1536 & 0.725 & 0.725 \\
      & $P_2 \oplus b_{\mathrm{face}}/P_0$
                                  & 3 & 1962 & 162  & 162  & 0.917 & 0.917 \\
      &                           & 4 & 4440 & 384  & 384  & 0.914 & 0.914 \\
\bottomrule
\end{tabular}
\caption{Assembled count at $\kappa/\mu = 10^4$. The bound is $1 - n_p/n_u$
and is undefined where $n_p > n_u$. Assembled zero-energy modes were zero in
every row. Script: \code{prototype/locking.jl}.}
\label{tab:assembled}
\end{table}

Three findings follow.

\paragraph{The constant-pressure soft modes do not persist after assembly.}
No spurious mode is present under either boundary condition at any mesh. Nor
are they present as low-energy modes: the smallest eigenvalue relative to the
deviatoric scale is $2.39\times10^{-3}$ for $P_0$ against $2.74\times10^{-3}$
for $P_1^{\mathrm{disc}}$ on the same mesh with the base alone fixed, and both
are the long-wavelength modes of that block. The three
modes per element of \S\ref{sec:firstresult} exist on each element and are element-local:
their traces do not match across faces.

What this test does not detect is the mode
\citet{Foulk.etal:2021} report, which is remedied there by a penalty and therefore has energy, small and
not zero. Those are the displacement modes in $\ker ~G_{P_0}$ whose
divergence is nonzero pointwise, which carry
shear energy alone in the element; whether they are spurious is the coercivity question of
\S\ref{sec:prediction}, which \S\ref{sec:plastic} measures, and not a count of
zeros. The composite tetrahedron of \citet{Foulk.etal:2021} is
built on a subdivided-linear displacement space rather than $P_2$, which is a
different kernel and is not modeled here; the soft mode of that element may be
specific to it. Nothing above contradicts their observation; it does not
reproduce it on a quadratic tetrahedron.

\paragraph{The unenriched pair locks.} With the boundary fixed,
$P_2/P_1^{\mathrm{disc}}$ retains 3--7\% of its deformation modes as $\kappa$
grows, against 58--63\% for $P_0$; with the base alone fixed, 37--41\%
against 80\%. Under either boundary condition it is the pair that
over-constrains, and locking is the more serious failure because the element
returns a stiff response and no spurious mode indicates the error.

\paragraph{The interior bubble raises the rank of $~G$ over a linear pressure.}
As Remark~\ref{rem:whyface} states, the interior bubble adds nothing to the
rank over $P_0$. Against the linear pressure modes its columns
are not zero: it raises $\operatorname{rank}~G$ from 957 to 1532 at $N = 4$,
and the interior-only pair has an isochoric fraction of $0.252$--$0.298$. The
full enrichment reaches $0.61$--$0.63$, equal to $P_0$'s to two digits and
without $P_0$'s element-level degeneracy; it attains the bound $1 - n_p/n_u$
to three digits and the interior-only pair to two, so every pressure mode is
resisted by a distinct displacement direction; and for the full pair $n_p - \operatorname{rank}~G = 1$ at every mesh,
the hydrostatic mode alone, where $P_2/P_0$ and the interior-only pair carry
four: the enrichment raises $\operatorname{rank}~G$ to $n_p - 1$.

\subsection{Stability, measured against a validated control}
\label{sec:beta}

The count of \S\ref{sec:infsup} is a necessary condition; Table~\ref{tab:beta}
shows that it is not sufficient.

Two properties of the computation allow meshes to $N = 8$. With $~K = \transpose{~P} ~L
\transpose{~L} ~P$ from a sparse Cholesky, $\transpose{q} ~G ~K^{-1}
\transpose{~G} q = \lVert ~L^{-1} (\transpose{~G} q)_{[p]} \rVert^2$, so
$\beta_h$ is the smallest nonzero singular value of $~W = ~L^{-1}
\transpose{~G} ~M^{-1/2}$ --- never forming $~G ~K^{-1} \transpose{~G}$, and so
never squaring the condition number of the quantity measured. And
the null dimension is fixed independently as $n_p - \operatorname{rank}~G$, so
the cut is taken by index rather than by selecting a threshold. The gap is
unambiguous: $O(10^{-15})$ below it against $O(10^{-2})$ above.

The sweep carries a positive control, the Taylor--Hood pair, whose stability
in three dimensions is established \citep{BoffiBrezziFortin:2013}: a valid
test must report it stable. A negative control alone cannot detect a test that
reports every pair unstable; the negative control, $P_1/P_0$, gives
$\beta_h = 0.333$, $0.229$, $0.165$, $0.129$ at $N = 2$ to $5$ with $45$,
$138$, $303$, $558$ unresisted pressure modes, the alternating-sign spurious
pressure mode of that pair (script \code{prototype/infsup.jl}).

\begin{table}[h]
\centering
\begin{tabular}{@{}lrrrrrr@{}}
\toprule
Pair & $h = 1/2$ & $1/4$ & $1/6$ & $1/8$ & null & local rate \\
\midrule
$P_2/P_1$ continuous (control) & 0.1734 & 0.2186 & 0.2214 & \textbf{0.2216} & 1 & $-0.003$ \\
\midrule
$P_2/P_0$                      & 0.1001 & 0.0755 & 0.0544 & 0.0418 & 4 & $0.97$ \\
$\oplus\, b_{\mathrm{int}}$    & 0.1015 & 0.0756 & 0.0544 & --- & 4 & $0.84$ \\
$\oplus\, b_{\mathrm{face}}$   & 0.4671 & 0.4228 & 0.4060 & \textbf{0.3967} & 1 & $0.076$ \\
\midrule
$P_2/P_1^{\mathrm{disc}}$      & 0.1132 & 0.0424 & 0.0277 & --- & $n_p - n_u$ & $1.04$ \\
$\oplus\, b_{\mathrm{int}}$    & 0.0782 & 0.0564 & 0.0406 & --- & 4 & $0.83$ \\
$\oplus\, b_{\mathrm{face}}$   & 0.0938 & 0.0701 & 0.0497 & --- & 2 & $0.88$ \\
$\oplus\, b_{\mathrm{int}}, b_{\mathrm{face}}$
                               & 0.2875 & 0.2962 & \textbf{0.2968} & --- & 1 & $-0.003$ \\
\bottomrule
\end{tabular}
\caption{$\beta_h$ over a Freudenthal mesh sequence, whole boundary fixed, ordered as
the construction of \citet[Example~8.7.2]{BoffiBrezziFortin:2013} runs. ``null''
is $n_p - \operatorname{rank}~G$, the pressure modes no displacement resists,
which must be one; ``local rate'' is $\ln(\beta_N/\beta_{N+1})/\ln((N+1)/N)$
at the last refinement, with $\beta_N$ the constant at refinement $N$, constant for a genuine power law and shrinking toward
zero for a sequence approaching a limit. Script: \code{prototype/beta.jl}.}
\label{tab:beta}
\end{table}

The control is reported stable: $\beta_h$ rises and settles on $0.2216$, flat
to three digits over the last three refinements, which validates the reading
of the remaining rows.

Every row of Table~\ref{tab:beta} is what \citet{BoffiBrezziFortin:2013}
predict. What the sweep establishes is that the scripts are sound and that the theory
applies to this mesh family; what it adds to the theory is the constants. The
rows are ordered as the construction of
\citet[Example~8.7.2]{BoffiBrezziFortin:2013} adds degrees of freedom.

\paragraph{$P_2$ does not control a constant pressure in three dimensions.}
$P_2/P_0$ decays at a local rate approaching one, although $n_p/n_u = 0.284$,
and carries four pressure null modes
at every mesh where one is allowed: the hydrostatic mode and three exact
spurious modes that no quadratic displacement resists. That is the statement
of \citet[\S8.7]{BoffiBrezziFortin:2013} quoted in \S\ref{sec:fe-spaces}, and
the three modes are the pressure modes that face degrees of freedom would
resist.

\paragraph{The interior bubble cannot supply it.} $P_2 \oplus
b_{\mathrm{int}}/P_0$ reproduces $P_2/P_0$ to three digits with the same four
null modes, for the reason in Remark~\ref{rem:whyface}. Over
$P_1^{\mathrm{disc}}$ the interior bubble changes $\beta_h$ from $0.0424$ to
$0.0564$ at $h = 1/4$ and from $0.0277$ to $0.0406$ at $h = 1/6$, and the
rate falls from $1.04$ to $0.83$, but the sequence decays and four null modes
remain. This pair has $n_p/n_u = 0.613$ (Table~\ref{tab:dofcount}) and is
unstable.

\paragraph{Face bubbles control the constant pressure.} $P_2 \oplus
b_{\mathrm{face}}/P_0$ drops the null dimension to one at every mesh and its
$\beta_h$ approaches from above with a local rate that falls from $0.16$ to
$0.08$ over the sequence, below the constant rate of a power law and above the
rate $-0.003$ of the control.
The seven refinements measured, of which Table~\ref{tab:beta} shows four,
cannot separate a limit near $0.35$ from a slow decay on their own, and the
row is reported as consistent with stable; its
stability rests on \citet[Example~8.7.2]{BoffiBrezziFortin:2013}, whose proof
constructs, from the face degrees of freedom, a displacement that matches the
divergence of any given field in the mean on every face.
Over $P_1^{\mathrm{disc}}$ the face bubbles alone leave one spurious linear
pressure mode per mesh and $\beta_h$ decays: controlling the linear modes is
the function of the interior bubble.

\paragraph{Both together are the three-dimensional Crouzeix--Raviart pair, and
it is stable.} $\beta_h$ rises monotonically to $0.2968$, flat to three digits
over the last three refinements, with a local rate of $-0.003$ --- the
same rate as the control's, and at a higher constant than the control's
$0.2216$. The pressure null space is one dimensional at every mesh. Each part
of the enrichment controls one part of the pressure space, and neither
controls the other's.

\begin{remark}[The distinct functions of the two enrichments]
\label{rem:whyface}
The interior bubble has zero pairing with a constant pressure: it vanishes on
the whole element boundary, so $\int_{\Omega_e}\!\Div ~b_{\mathrm{int}} = 0$
and its columns of $~G$ against a $P_0$ mode are identically zero. It pairs
only with the linear pressure modes, which is why it changes
$P_2/P_1^{\mathrm{disc}}$ and changes the $P_2/P_0$ constant by at most
$0.0014$ (the $P_2 \oplus b_{\mathrm{int}}/P_0$ row of Table~\ref{tab:beta},
with the same null dimension): upward, because enlarging $U_h$ cannot lower
$\beta_h$, and by at most $0.0014$, because the bubble's columns of $~G$
against $P_0$ are zero.
\code{checks.jl} reads those columns directly and confirms it. The face bubbles are what supply a displacement response to
the pressure jump across a face, which is the mode a discontinuous pressure
needs resisted.
\end{remark}

$P_2/P_0$ and $P_2 \oplus b_{\mathrm{int}}/P_1^{\mathrm{disc}}$ have
$n_p/n_u = 0.284$ and $0.613$ and are unstable: a ratio above one excludes a
pair, and a ratio below one does not admit it.

\begin{remark}[Mesh dependence]
\label{rem:structured}
These results are for the Freudenthal subdivision, a structured family, and
structured meshes are where pair counterexamples characteristically arise:
the alternating-sign spurious pressure mode of $P_1/P_0$ has its largest
amplitude on them. A pair can degenerate on a structured mesh and remain well
behaved on a general one. On this family, $P_2/P_0$,
$P_2/P_1^{\mathrm{disc}}$ and $P_2 \oplus b_{\mathrm{int}}/P_1^{\mathrm{disc}}$
decay; confirming that as a property of the pairs rather than of the mesh
needs a sequence of unstructured meshes, which this study does not have.
\S\ref{sec:deformed} gives one measurement in the opposite sense: on the
non-affine images of this mesh the $P_2/P_0$ constant is smaller. The affirmative result
needs no such confirmation, since \citet[Example~8.7.2]{BoffiBrezziFortin:2013}
establish it on any shape-regular family; what the sweep adds is the constant.
\end{remark}

\subsection{Stability on a deformed configuration}
\label{sec:deformed}

Every measurement above is at $\F = \Id$. At a finite deformation the
pairing is $\int \varpi\,\mathrm{div}_{\bx}\testd \diff V$ (\S\ref{sec:fe-spaces}),
the divergence pairing on the deformed configuration with the pressure scaled
by $1/J$, and whether $\beta_h$ stays bounded below under a large non-affine
deformation is what the regime of \S\ref{sec:problem} requires. The operators
\eqref{eq:infsup-ops} are assembled on the isoparametric image of the
quadratic mesh under a prescribed map, with the derivatives, the divergence
and the volume taken on the image and the pressure basis on the reference
cell as the element uses it, and with the boundary conditions of the
reference cube; the scaling by $1/J$ changes $\beta_h$ by at most the ratio
of extreme Jacobians and is not applied. Five maps of the unit cube, about
its center: a twist about the vertical axis by $90^\circ$ and by $180^\circ$
over the height, isochoric and non-affine, with shear strains near $2$ at the
corners of the latter; a parabolic shear $x_1 = X_1 + \tfrac12 X_3^2$,
isochoric; a radial inflation with $J$ from $1$ to $2.3$; and an affine
compression to $J = \tfrac12$ as a control, since an affine map changes
$\beta_h$ only through the element aspect ratio. Table~\ref{tab:deformed}
gives the result at $N = 5$ with the local rate over the last refinement.
Script: \code{prototype/deformed.jl}.

\begin{table}[tbp]
\centering
\caption{$\beta_h$ at $N = 5$ on deformed configurations, whole boundary
fixed, with the local rate over $N = 4 \to 5$ in parentheses for the
Crouzeix--Raviart pair; the null dimension is one on every row for the three
left-hand pairs. The $P_2/P_0$ column gives its null dimension and
$\beta_h$: its three exact spurious modes exist only on the affine mesh.}
\label{tab:deformed}
\begin{tabular}{@{}llrrrl@{}}
\toprule
map & $J$ & $P_2/P_1$ cont. & $P_2 \oplus b_{\mathrm{face}}/P_0$
    & $P_2 \oplus b_{\mathrm{int}} \oplus b_{\mathrm{face}}/P_1^{\mathrm{disc}}$
    & $P_2/P_0$ (null, $\beta_h$) \\
\midrule
reference          & $1$          & 0.221 & 0.413 & 0.297 ($-0.003$) & 4, 0.063 \\
twist $90^\circ$   & $0.99$--$1.01$ & 0.189 & 0.385 & 0.204 ($0.18$)  & 1, 0.027 \\
twist $180^\circ$  & $0.98$--$1.05$ & 0.137 & 0.345 & 0.116 ($0.42$)  & 1, 0.051 \\
parabolic shear    & $1$          & 0.176 & 0.336 & 0.183 ($0.05$)   & 1, 0.005 \\
inflation          & $1.0$--$2.3$   & 0.192 & 0.368 & 0.276 ($0.01$)   & 1, 0.004 \\
compression        & $0.5$        & 0.192 & 0.333 & 0.213 ($0.00$)   & 4, 0.038 \\
\bottomrule
\end{tabular}
\end{table}

Three findings follow.

\paragraph{$\beta_h$ is not uniform in the deformation.} It decreases with
the local shear of the map for every pair, and most for the
Crouzeix--Raviart pair: from $0.297$ to $0.204$ under the $90^\circ$ twist,
$0.116$ under the $180^\circ$ twist and $0.183$ under the parabolic shear,
against $0.221 \to 0.189$, $0.137$, $0.176$ for Taylor--Hood and
$0.413 \to 0.385$, $0.345$, $0.336$ for the face-bubbled constant-pressure
pair. The inflation, which changes $J$ by a factor of $2.3$ but shears
little, leaves the constant at $0.276$. The affine compression gives $0.213$,
flat to four digits from $N = 3$: the aspect-ratio effect alone.
The dependence is on the shear of the deformation and not on its volume
change.

\paragraph{At a fixed deformation the sequence in $h$ is consistent with a
positive limit.} The local rate of the Crouzeix--Raviart pair falls with
$N$ on every non-affine map: $0.28$, $0.22$, $0.18$ over $N = 2$ to $5$ for
the $90^\circ$ twist, $0.67$, $0.53$, $0.42$ for the $180^\circ$ twist, and
$0.05$ for the parabolic shear (through $N = 5$). A power-law decay would hold its rate. The
$180^\circ$ row has not flattened by $N = 5$, and its shear strain of about
$2$ at the corners is beyond what a plastic zone of the problems listed in
\S\ref{sec:verify} reaches; the finding is that the constant degrades with the strain, not that
it vanishes. At $N = 6$ the pair gives $0.1975$, $0.1088$ and $0.182$ on
the three non-affine maps, local rates of $0.17$, $0.35$ and $0.04$, each
below the rate before it.

\paragraph{The exact spurious modes of $P_2/P_0$ belong to the affine mesh.}
Under every non-affine map the pressure null dimension of $P_2/P_0$ falls
from four to one, the three modes are resisted, and $\beta_h$ of $P_2/P_0$
is smaller on every non-affine image, $0.004$ to $0.05$, than on the reference
mesh, $0.063$.

\subsection{Coercivity under a plastic tangent}
\label{sec:plastic}

The claim of \S\ref{sec:prediction} is measured as the constant it names. With
$~Z$ a basis of $\ker ~G_h$,
\begin{equation}
  \label{eq:coercivity}
  \lambda_{\min} = \min_{q}
    \frac{\transpose{q}\,\transpose{~Z} ~K^{\mathrm{ep}} ~Z\, q}
         {\transpose{q}\,\transpose{~Z} ~K^{H^1} ~Z\, q},
  \qquad
  ~K^{\mathrm{ep}}_{ij} = \frac{1}{2\mu}\int \eps_i : \mathbb{C}^{\mathrm{ep}} : \eps_j \diff V,
\end{equation}
is the smallest generalized eigenvalue of the deviatoric form under the
tangent \eqref{eq:cep} against the $H^1$ seminorm, restricted to the
discretely isochoric subspace. It is dimensionless, it is taken on the exact
kernel rather than through a penalty, and it is the constant of the second
hypothesis; \eqref{eq:floor} says what to expect of it, at least
$(1-\eta)/2$ for every pair when $\eta < 1$ and a decay to zero for every
pair when $\eta = 1$. The quantity that carries the claim is the dilatation
of the eigenmodes,
\begin{equation}
  \label{eq:rv}
  r_v(~v) := \frac{\int (\Div ~v)^2 \diff V}{\int \lvert\Grad ~v\rvert^2 \diff V},
\end{equation}
zero for an isochoric field. The flow direction is uniform, $~n =
\sym(~e_1 \otimes ~e_2)/\sqrt{2}$ or $~n = \dev(~e_3 \otimes ~e_3)/\norm{\cdot}$,
which is the most unfavorable choice: the continuum operator then admits shear
bands of any width, a decaying eigenvalue is expected of every pair, and only
the dilatation of the mode says whether the band is physical. The tangent is
taken at $\eta = 1$ and at $\eta = 0.9$, the boundary is fixed everywhere,
and each pair is measured elastically ($\eta = 0$) on the same kernel.
Script: \code{prototype/plastic.jl}.

\begin{table}[tbp]
\centering
\caption{Lowest mode of the deviatoric form on $\ker ~G_h$ under the tangent
\eqref{eq:cep} at $\eta = 1$, Freudenthal cube with the whole boundary fixed: the eigenvalue
$\lambda_1$ of \eqref{eq:coercivity} and the dilatation $r_v$ of its mode,
\eqref{eq:rv}, for a shear and a uniaxial flow direction. The elastic lowest
mode has $\lambda_1 = 0.500$ and $r_v = 0$ on every row, as the identity
behind \eqref{eq:floor} requires. At $\eta = 0.9$ every $\lambda_1$ exceeds
the lower bound $0.05$ of \eqref{eq:floor} and the dilatations agree with those below to two digits.
Script: \code{prototype/plastic.jl}.}
\label{tab:plastic}
\begin{tabular}{@{}lrrrrr@{}}
\toprule
& & \multicolumn{2}{c}{$~n = \sym(~e_1 \otimes ~e_2)/\sqrt2$}
  & \multicolumn{2}{c}{$~n = \dev(~e_3 \otimes ~e_3)/\norm{\cdot}$} \\
\cmidrule(lr){3-4}\cmidrule(lr){5-6}
Pair & $N$ & $\lambda_1$ & $r_v$ & $\lambda_1$ & $r_v$ \\
\midrule
$P_2/P_0$                       & 3 & 0.109  & 0.065  & 0.117 & 0.208 \\
                                & 4 & 0.067  & 0.034  & 0.106 & 0.217 \\
                                & 5 & 0.047  & 0.029  & 0.101 & \textbf{0.220} \\
$P_2 \oplus b_{\mathrm{face}}/P_0$
                                & 3 & 0.044  & 0.044  & 0.082 & 0.274 \\
                                & 4 & 0.026  & 0.028  & 0.074 & 0.268 \\
                                & 5 & 0.017  & 0.019  & 0.071 & \textbf{0.261} \\
$P_2 \oplus b_{\mathrm{int}} \oplus b_{\mathrm{face}}/P_1^{\mathrm{disc}}$
                                & 3 & 0.074  & 0.036  & 0.147 & 0.018 \\
                                & 4 & 0.049  & 0.019  & 0.139 & 0.0074 \\
                                & 5 & 0.032  & 0.0066 & 0.134 & \textbf{0.0042} \\
$P_2/P_1$ continuous (control)  & 3 & 0.056  & 0.051  & 0.053 & 0.861 \\
                                & 4 & 0.030  & 0.029  & 0.026 & 0.929 \\
                                & 5 & 0.019  & 0.020  & 0.016 & \textbf{0.956} \\
\bottomrule
\end{tabular}
\end{table}

\paragraph{The scripts are checked by an identity.} By the identity behind
\eqref{eq:floor}, the elastic form on any kernel satisfies
$\lambda = \tfrac12 + \tfrac16 r_v$, so its lowest mode has $r_v = 0$ and
$\lambda_1 = 0.5$ exactly. Every row of the sweep reproduces this; at $N = 2$
for $P_2/P_0$, where the kernel has no isochoric member, $\lambda_1 = 0.5024$
with $r_v = 0.0146$, as the identity requires. The identity is also the check
on the energy metric of the scripts.

\paragraph{The constant does not separate pairs, as \eqref{eq:floor}
predicts.} Under the shear direction $\lambda_1$ decays for every pair, the
control included, at rates between $h^{1.5}$ and $h^2$: the discrete spaces
resolve narrower shear bands, which the continuum operator admits with zero energy.
Under the uniaxial direction no shear band has $\eps \propto ~n$, and
$\lambda_1$ converges for the three element-local pairs, to $0.10$ for
$P_2/P_0$, $0.071$ for $P_2 \oplus b_{\mathrm{face}}/P_0$ and $0.13$ for the
Crouzeix--Raviart pair, while it decays for Taylor--Hood. At $\eta = 0.9$
every $\lambda_1$ lies above the lower bound $0.05$ of \eqref{eq:floor}, the smallest being the
control's $0.067$.

\paragraph{The dilatation of the soft modes separates them, in the direction
claimed.} Under the uniaxial direction the lowest modes of both
constant-pressure pairs carry $r_v = 0.20$--$0.27$ at every refinement, with
no decay. These are the modes of \S\ref{sec:prediction}: the element assigns
them no energy for a dilatation the continuum resists with $\kappa$, and they
are the softest modes the element has. The lowest modes of the
Crouzeix--Raviart pair carry $0.040$, $0.018$, $0.0074$ and $0.0042$ over
$N = 2$ to $5$, decaying as $h^2$, the rate at which
\citet[Prop.~5.1.3]{BoffiBrezziFortin:2013} makes $\ker ~G_h$ approximate the
isochoric space; its soft modes are physical. Under the shear direction the
lowest modes of every pair are shear bands whose dilatation decays, the
Crouzeix--Raviart pair's three to four times smaller than the others' at
$N = 5$. Counted over the twenty lowest modes, with a mode called spurious when
$r_v > 0.1$ (the stable pair's soft modes lie below $0.02$ from $N = 4$ on,
the spurious ones at $0.2$--$0.4$): under the uniaxial direction at $N = 5$
all twenty of each constant-pressure pair's and of Taylor--Hood's are
spurious, and none of the Crouzeix--Raviart pair's at any $N$; under the
shear direction the twenty lowest of every element-local pair are bands at
$N = 5$. The spurious modes of a constant pressure are therefore present in
both directions but are the lowest modes only where the flow admits no band;
the next paragraph therefore counts them among the twenty lowest in a zone
with a varying flow direction.

\paragraph{A confined zone with a varying flow direction.} The flow
direction and the plastic zone are taken from a prescribed displacement
field: $~n = \dev\eps(~u)/\norm{\dev\eps(~u)}$ at each quadrature point and
$\eta = 1$ where $\lvert\dev\eps(~u)\rvert$ exceeds $0.35$ of its peak,
$\eta = 0$ elsewhere. Two fields: an indentation of the top face,
$u_3 = -\exp(-d^2/w^2)\,X_3$ with $d$ the distance from the vertical axis and
$w = \tfrac14$, whose zone is $51\%$ of the volume with elastic
surroundings; and the displacement of the $90^\circ$ twist of
\S\ref{sec:deformed}, whose zone is $97\%$. Both are run with the whole
boundary fixed and with the base alone fixed. Table~\ref{tab:zone} counts
the spurious modes among the twenty lowest. Script:
\code{prototype/plastic\_zone.jl}.

\begin{table}[tbp]
\centering
\caption{Spurious modes among the twenty lowest of the plastic form on
$\ker ~G_h$ ($r_v > 0.1$) in a confined plastic zone with a varying flow
direction, at $N = 3$, $4$, $5$.}
\label{tab:zone}
\begin{tabular}{@{}llrrrr@{}}
\toprule
field & boundary & $P_2/P_0$ & $P_2 \oplus b_{\mathrm{face}}/P_0$
      & $P_2 \oplus b_{\mathrm{int}} \oplus b_{\mathrm{face}}/P_1^{\mathrm{disc}}$
      & $P_2/P_1$ cont. \\
\midrule
indentation & whole fixed & 16, 15, 7  & 18, 17, 5  & 4, \textbf{0, 0} & 20, 18, 17 \\
indentation & base fixed & 11, 1, 0   & 13, 5, 0   & \textbf{0, 0, 0} & 14, 7, 5 \\
twist       & whole fixed & 14, 18, 15 & 18, 13, 8  & 1, 4, \textbf{0} & 20, 20, 20 \\
twist       & base fixed & 5, 2, 0    & 13, 7, 4   & \textbf{0, 0, 0} & 12, 13, 12 \\
\bottomrule
\end{tabular}
\end{table}

With the whole boundary fixed the constant-pressure pairs retain between
five and fifteen spurious modes among their twenty softest at $N = 5$, and
under the twist the softest mode of $P_2/P_0$ is spurious at every $N$ up to
$4$ ($r_v = 0.22$--$0.31$) and its third at $N = 5$ ($0.20$). The
Crouzeix--Raviart pair has none from $N = 4$ on under the indentation and
none at $N = 5$ under the twist, and none at any $N$ with the base alone
fixed. Taylor--Hood keeps seventeen to twenty, and under the twist its
softest mode's dilatation grows with refinement, $0.29$, $0.37$, $0.66$,
$0.74$. With the base alone fixed the softest modes of every pair are
long-wavelength motions of the block and the counts fall for the
constant-pressure pairs too, to zero at $N = 5$ for $P_2/P_0$ and four for
the face-bubbled pair, so the spurious modes are then above the twentieth;
they are not absent, as the whole-boundary rows show. Whether a nonlinear
solve excites them is the open measurement of \S\ref{sec:verify}.

\paragraph{Taylor--Hood is the worst pair on this measure.} Its pressure is
continuous, so its kernel is divergence-free only against continuous linear
test functions, and under the uniaxial direction its lowest modes are
dilatation almost entirely: $r_v = 0.86$, $0.93$, $0.96$ at $N = 3$, $4$, $5$,
with $\lambda_1$ decaying to $0.016$. Those modes have no counterpart in the
continuum, which resists them with $\kappa$. Taylor--Hood is therefore an
alternative for elasticity only (\S\ref{sec:risks}).


\subsection{Quadrature rule and basis}
\label{sec:basis}

Two implementation choices set the number of operations per element and its
explicit stable step, and neither is decided by the stability results: the
quadrature rule, and the basis that carries the enriched space. Both are
measured on one element. Scripts:
\code{prototype/quadrature.jl} and \code{prototype/basis.jl}.

\paragraph{Quadrature.} The scripts integrate the enriched element to degree
$6$ so that the linear measurements are exact. For the element the criterion
is weaker: a rule is admissible when the condensed stiffness
$~K_e = ~K\supdev + \kappa\transpose{~G}~M^{-1}~G$ has exactly six
zero-energy modes, and its effect is the softening it introduces, measured as
the extreme generalized eigenvalues of $~K_e(\text{rule})$ against
$~K_e(\text{exact})$ on the complement of the rigid modes, which are $1$ for
an exact rule. Table~\ref{tab:quadrature} gives both for the rules
considered, with a degree-11 rule of $216$ points as the reference. The two
rules supplied by \code{ReferenceFiniteElements} (RFE), the Julia package that
provides the reference elements and their quadrature, four and five points,
leave
$21$ zero modes and are inadmissible. The eight-point conical rule (a
tensor-product Gauss rule mapped onto the tetrahedron) is rank sufficient and softens the element to $0.03$--$0.10$ of its stiffness in
the worst direction, which is also inadmissible. Two degree-5 rules are
admissible: the $14$-point symmetric rule of \citet{Keast:1986}, which
integrates every monomial through degree $5$ to $4\times10^{-16}$ and
over-stiffens the interior bubble's own block by $35\%$ since its
gradient-squared is of degree $6$, and the $27$-point conical rule, which
under-stiffens the same block by $9$--$12\%$. Assembled, both give the
Crouzeix--Raviart pair's constant within $1\%$ ($0.293$ and $0.298$ against
$0.297$ at $N = 5$), a one-dimensional pressure null space, no spurious mode
among the twenty lowest under the uniaxial plastic tangent, and the same
dilatation of the lowest mode ($0.0073$, $0.0074$ and $0.0074$). The
$14$-point rule satisfies the criteria --- six zero-energy modes, positive
weights, $4.6$ times fewer points than the $64$-point degree-7 rule --- and
is used, with a stiffer interior bubble.

\begin{table}[tbp]
\centering
\caption{Quadrature rules on the enriched element: zero-energy modes of the
condensed stiffness (six is correct) and of the deviatoric stiffness (ten:
six rigid modes and the four volumetric directions of
\S\ref{sec:firstresult}), and the extreme softening
factors against a $216$-point rule, on the reference tetrahedron and on a
Kuhn tetrahedron of the Freudenthal mesh. $\kappa/\mu = 10^4$. Script:
\code{prototype/quadrature.jl}.}
\label{tab:quadrature}
\begin{tabular}{@{}lrrrrrrr@{}}
\toprule
& & \multicolumn{3}{c}{reference} & \multicolumn{3}{c}{Kuhn} \\
\cmidrule(lr){3-5}\cmidrule(lr){6-8}
rule & points & zero($~K_e$) & min & max & zero($~K_e$) & min & max \\
\midrule
RFE degree 2            & 4   & 21 & 0     & 1.30 & 21 & 0     & 1.30 \\
RFE degree 3            & 5   & 21 & $<0$  & 1.13 & 21 & $<0$  & 1.12 \\
conical degree 3        & 8   & 6  & 0.10  & 1.14 & 6  & 0.03  & 1.21 \\
Keast degree 5          & 14  & 6  & 1.00  & 1.35 & 6  & 1.00  & 1.35 \\
conical degree 5        & 27  & 6  & 0.91  & 1.00 & 6  & 0.88  & 1.00 \\
conical degree 7        & 64  & 6  & 1.00  & 1.00 & 6  & 1.00  & 1.00 \\
\bottomrule
\end{tabular}
\end{table}

\paragraph{Basis.} The space $P_2 \oplus b_{\mathrm{int}} \oplus
b_{\mathrm{face}}$ can be carried by the hierarchical basis of the scripts,
ten Lagrange functions and five bubbles, or by the nodal basis with unit value
at one of fifteen nodes, the \code{TETRA15} arrangement: vertices, edge
midpoints, face centroids, centroid. The space is the same, so every
stiffness and stability result is the same; what differs is conditioning,
lumping, and the coefficients of a rigid translation. Table~\ref{tab:basis}
gives the measurements at $\kappa/\mu = 2$. The nodal basis conditions the
consistent mass eight times better ($94$ against $778$) and the stiffness
better ($559$ against $797$). Under HRZ lumping the hierarchical basis loses
$51\%$ of the momentum of a rigid translation unless the normalization runs
over the Lagrange functions alone (\S\ref{sec:dynamics}); the nodal basis conserves it without
qualification, and its lumped masses are positive on every element
measured. The explicit stable step from the element bound, as the Courant
number $c_p\Delta t/h$, is $0.476$ for \code{TETRA4}, $0.197$ for
\code{TETRA10} under HRZ, $0.145$ for the enriched element in the nodal
basis under HRZ, and $0.114$ in the hierarchical basis under Lagrange-only
HRZ, with the same ordering on the distorted and Kuhn elements. The nodal
basis is used: every unknown is attached to a node, so no unknown of a new
kind is introduced; its mass lumps correctly by the standard construction;
and it takes a $20$--$27\%$ larger stable step.

\begin{table}[tbp]
\centering
\caption{Basis choice on the reference tetrahedron, $\kappa/\mu = 2$,
$\mu = \rho = 1$: condition numbers of the condensed stiffness on the
complement of the rigid modes and of the consistent mass, the smallest
lumped mass over $\rho V$, the momentum of a unit translation under the
lumped mass over $\rho V$, and the Courant number $c_p\Delta t/h$ from the
element bound. Script: \code{prototype/basis.jl}.}
\label{tab:basis}
\begin{tabular}{@{}lrrrrr@{}}
\toprule
basis, lumping & cond $~K_e$ & cond $~M$ & min $m/\rho V$ & momentum & Courant \\
\midrule
hierarchical, consistent            & 797 & 778 & ---    & ---   & 0.092 \\
hierarchical, HRZ over all 15       & --- & --- & 0.014  & 0.486 & 0.097 \\
hierarchical, HRZ over Lagrange 10  & --- & --- & 0.028  & 1.000 & 0.114 \\
nodal, consistent                   & 559 & 94  & ---    & ---   & 0.092 \\
nodal, HRZ                          & --- & --- & 0.021  & 1.000 & 0.145 \\
\code{TETRA10}, HRZ                 & --- & --- & 0.028  & 1.000 & 0.197 \\
\code{TETRA10}, row sum             & --- & --- & $-0.050$ & 1.000 & --- \\
\code{TETRA4}, row sum              & --- & --- & 0.250  & 1.000 & 0.476 \\
\bottomrule
\end{tabular}
\end{table}

\section{Summary of the evidence and open measurements}
\label{sec:verify}

The formulation rests on five claims. For each, the section that establishes
it and the number that carries it are given; the measurements that remain are
listed at the end.

\paragraph{The pair is inf-sup stable in three dimensions.} \S\ref{sec:beta}:
the Crouzeix--Raviart pair has $\beta_h = 0.2968$ at a local rate of $-0.003$,
above the Taylor--Hood control's $0.2216$, with a one-dimensional pressure
null space at every mesh; the constant-pressure pair with face bubbles is
stable with a constant decreasing toward a limit near $0.35$, and the five
other element-local pairs decay. The face bubbles are the necessary part of
the enrichment; the interior bubble alone is
not sufficient (Remark~\ref{rem:whyface}). Uniformity in $\kappa/\mu$ follows
from stability of the incompressible limit
\citep[\S5.5.2, \S8.12.1, Remark~8.12.1]{BoffiBrezziFortin:2013}. On deformed
configurations (\S\ref{sec:deformed}) the constant decreases with the shear
of the deformation, to $0.204$ at a $90^\circ$ twist and $0.116$ at
$180^\circ$, and by $7\%$ under a volume change of $2.3$, and at each fixed
deformation the sequence in $h$ is consistent with a positive limit; the
constant is not uniform in the deformation.

\paragraph{No spurious soft modes.} In elasticity the second hypothesis of
Brezzi's theorem holds for every conforming pair (\S\ref{sec:prediction}). The
three exact zero-energy modes per element of a constant pressure
(\S\ref{sec:firstresult}) do not persist after assembly
(\S\ref{sec:assembled}). Under a plastic tangent (\S\ref{sec:plastic}) the
coercivity constant is bounded below by \eqref{eq:floor} for every pair, and
the dilatation of the soft modes separates them: $0.20$--$0.27$ without
decay for the constant-pressure pairs, $h^2$ decay for the Crouzeix--Raviart
pair, and in a confined zone with a varying flow direction no spurious mode
among the twenty softest of the Crouzeix--Raviart pair at $N = 5$ against
five to fifteen for the constant-pressure pairs.

\paragraph{No volumetric locking.} \S\ref{sec:assembled}: the discretely
isochoric fraction $\dim\ker ~G / n_u$ is $0.61$--$0.63$ for the
Crouzeix--Raviart pair, at the bound $1 - n_p/n_u$, against
$0.03$--$0.07$ for $P_2/P_1^{\mathrm{disc}}$, which locks.

\paragraph{Element-local elimination.} \S\ref{sec:local}: the pressure is
discontinuous, so its elimination is element by element and the tangent has
the sparsity of the displacement mesh; the face bubbles add global
displacement unknowns, $2.6$ times those of $P_2/P_0$ at $N = 12$
(\S\ref{sec:infsup}).

\paragraph{Quadrature and basis.} \S\ref{sec:basis}: the $14$-point degree-5
rule keeps six zero-energy modes and the assembled constant within $1\%$; the
nodal basis conserves the momentum of a rigid translation under HRZ lumping
and gives a Courant number of $0.145$ against $0.476$ for \code{TETRA4}.

\paragraph{Open measurements.} Three remain. An unstructured mesh sequence
for the negative stability results, which are established here on the
Freudenthal family only (Remark~\ref{rem:structured}). The smoothness of the
equivalent plastic strain in a nonlinear solve, on the benchmark set on which
the composite tetrahedron performs poorly \citep{Foulk.etal:2021}: punch
indentation, a nearly incompressible Cook's membrane in three dimensions (a
tapered panel clamped on one edge and sheared on the opposite one),
large-rotation twist of a bar, and a weak shock in a material of high
$\kappa/\mu$; this is to be measured with the Simo--Hughes $J_2$ model
implemented in Norma and Carina rather than with the additive logarithmic
plasticity of \S\ref{sec:plasticity} (Remark~\ref{rem:whichj2}), against the
composite tetrahedron on the same meshes. And the sparsity of the assembled
tangent in an implementation, which \S\ref{sec:local} establishes by
construction.

\section{Limitations}
\label{sec:risks}

This section states what the formulation requires, what it does not
establish, and what its parts are.

\paragraph{The number of unknowns.} The face bubbles of the Crouzeix--Raviart
pair are shared between neighbors and do not condense: at $N = 12$ the
enriched space carries $96{,}117$ global displacement degrees of freedom
against $36{,}501$ for $P_2/P_0$ and for the composite tetrahedron on the
same mesh, a factor of $2.6$, and $2.5$ times Taylor--Hood's $38{,}698$
(\S\ref{sec:infsup}). Whether the accuracy of a stable pair on coarser meshes
compensates for that factor is not established; for elasticity that is the
comparison with Taylor--Hood, whose pressure is global and whose displacement
space is $2.5$ times smaller, and for plasticity Taylor--Hood's softest modes
are dilatation almost entirely (\S\ref{sec:plastic}), so the comparison there
is with the composite tetrahedron.

\paragraph{Operations per element.} With the $14$-point rule of
\S\ref{sec:basis} the enriched element evaluates the constitutive law at $14$
points on $45$ displacement degrees of freedom, against $4$ or $5$ points on
$30$ for \code{TETRA10} and one on $12$ for \code{TETRA4}. The number of
elements a quadratic tetrahedron needs for a given accuracy in this problem
class is not measured.

\paragraph{The explicit stable step.} HRZ lumping in the nodal basis gives
positive masses and conserves momentum, and the stable step is $3.3$ times
smaller than \code{TETRA4}'s on the same element, of which a factor of $2.4$
is the quadratic element and $1.4$ the enrichment (\S\ref{sec:basis}). The comparison at equal accuracy, on the coarser mesh a
quadratic element allows, is not measured; if weak shocks are run with the
implicit integrator the factor does not apply.

\paragraph{The material.} \S\ref{sec:matreq} requires an exact
volumetric--isochoric split. Of the seven materials implemented in Norma,
five satisfy it; linear elasticity and Saint-Venant--Kirchhoff do not, and
are not finite-deformation models in the sense of \S\ref{sec:problem}. The
three with a volumetric energy that is not quadratic require one projection
more (Remark~\ref{rem:nonquadratic}). Anisotropic or non-separable response
is outside the split and would need the five-field variant of
\S\ref{sec:fivefield} or a different treatment. The $\bar F$ family of
elements carries no such requirement.

\paragraph{The parts.} Mixed volumetric treatment
\citep{SimoTaylorPister:1985}, logarithmic-strain plasticity
\citep{Miehe.Apel.Lambrecht:2002} and displacement--pressure pairs satisfying
\eqref{eq:infsup} uniformly \citep{BoffiBrezziFortin:2013} are established.
The contribution is their combination: an element-local pair satisfying
\eqref{eq:infsup} uniformly, a projected volumetric strain, and a deviatoric
law evaluated pointwise, which removes the stabilization that the
composite-tetrahedron and variational-multiscale methods carry. The evidence
is that of \S\ref{sec:results}; the nonlinear confirmation is the open
measurement of \S\ref{sec:verify}. Methods with unknowns on element faces and
element interiors coupled only through the faces \citep{AbbasErnPignet:2019}
are the alternative should that measurement not confirm it.

\bibliographystyle{plainnat}
\bibliography{references}

\end{document}
